% SPDX-FileCopyrightText: © 2026 Idiap Research Institute <contact@idiap.ch>
%
% SPDX-FileContributor: Andre Freitas <andre.freitas@idiap.ch>
% Neuro-symbolic AI Group
%
% SPDX-FileContributor: Danilo Gusicuma <danilo.gusicuma@idiap.ch>
% Neuro-symbolic AI Group
%
% SPDX-License-Identifier: MIT
\clearpage
\section{Workflow mechanisms at Prism occurrence level}
\label{app:granular-prism-workflows}
A motif identifies a recurring arrangement of reasoning. A Prism-level description must also say which occurrence performs each move, which typed artifact it consumes, which relation connects it to another occurrence, and what information survives their composition. The following descriptions cover all nineteen canonical workflows at this level. They report the native declarations and distinguish the original workflow's routing policy from the native return boundary. Where the declaration exposes an unresolved projection, the native representation gap remains visible. The underlying skill procedures are described separately: their conceptual instructions can be richer than that projection, without establishing its implementation.

\paragraph{Reading the representation.}
Each subsection introduces local aliases for nominal artifact types; aliases reset between workflows. An operation $f:X\to\mathsf{Gen}[Y]$ produces a generated artifact with payload type $Y$, where $\mathsf{Gen}[Y]=\mathtt{Generated}[Y]$. The explicit projection $z.\mathtt{value}$ refers to a successfully obtained payload, not a proof of its proposition or a way to eliminate a model failure. Named records display joins: for example, a calibration input containing both a claim and its assessment preserves two distinct dependencies. Parallel occurrences can share an input without establishing independence of their scientific evidence.

A native occurrence has the schematic form
\begin{quote}\small\ttfamily
[result: Operation(input)] by Relation
\end{quote}
The binding names the resulting generated artifact; the operation gives the conceptual transformation; the relation records the license connecting source and target. At the declaration level its shape is
\[
 \Gamma\vdash x:X,\qquad\Gamma\vdash y:\mathsf{Gen}[Y],\qquad
 \rho(x,y):\mathsf{Certificate}_{\rho}[X,Y].
\]
This notation summarizes a declared relation contract. It does not assert that a certificate is a proof that $y$ is true or that all scientific preconditions of the transition hold. The claim-indexed refinement in Appendix~\ref{app:dependent-workflow} makes a stronger association explicit; it is not inferred from these nominal signatures.

\paragraph{Granularity and limits.}
Each occurrence uses seven compact fields: \textbf{Purpose}, \textbf{Input},
\textbf{Method}, \textbf{Check}, \textbf{Output}, \textbf{Next} and
\textbf{Example}. Input and Output name the carried scientific content;
Next records continuation or stopping separately. Numbering provides a reading order;
the typed dependencies determine execution constraints, including parallel branches.
The case remains continuous across steps, and a changed-input contrast explains
what would alter its conclusion. Source-prescribed procedures describe what the
linked skills ask for, not an implemented algorithm inferred from a native signature.

Every declared occurrence is described below, including projection holes and repeated calibrations with different responsibilities. The added procedure explanations draw on selected linked skill instructions and workflow contracts; a binding list alone establishes neither execution order nor realization of those instructions. A source-policy route to a later workflow is not automatically a typed composition: an additional artifact or adapter may be required. The examples are constructed scientific cases, not reported executions. Exact source names, relation mappings and source hashes accompany the editable paper assets; the prose explains their conceptual significance.

% SPDX-FileCopyrightText: © 2026 Idiap Research Institute <contact@idiap.ch>
%
% SPDX-FileContributor: Andre Freitas <andre.freitas@idiap.ch>
% Neuro-symbolic AI Group
%
% SPDX-FileContributor: Danilo Gusicuma <danilo.gusicuma@idiap.ch>
% Neuro-symbolic AI Group
%
% SPDX-License-Identifier: MIT
\subsection{Abductive proof work: making a proposed route answerable}
\label{app:granular-abducer-work-loop}
An active proof gap is an obligation, not yet a route for meeting it. The four
occurrences make four different commitments. Route formation ($a$) proposes a
hypothesis that could close the gap. Claim derivation ($d$) takes that hypothesis
as a premise and exposes a consequence. Refutation ($r$) challenges that exact
consequence. Routing ($e$) retains both the consequence and its assessment to
choose what work is justified next. This retained join matters: an adverse test
without the claim cannot say which part of the proposed route failed.

\paragraph{How the search is constrained.}
The linked abductive skill first attempts deduction to locate the exact missing
support, then constructs competing small sets of additional premises. It compares
explanatory coverage, compatibility with accepted knowledge, simplicity, and
novelty risk, and attaches a falsifier to each set. Thus the search space is not
``any plausible explanation'': it is alternative repairs of a named entailment
gap. The linked outer work procedure alternates discovery and proof attempts,
reassesses end-to-end plausibility, and permits persistence only when failure
sharpens a premise, object, regime, or bridge. Repeating the same blocker instead
requires backtracking or a pivot. Its machinery assessment asks which method
actually controls the bottleneck, excluding a prestigious analogy with no
transfer obligation. These are prescribed skill procedures; the four native
occurrences do not implement that persistent controller.

Here $G,H,C,R,S$ alias the native gap, route hypothesis, derived claim,
refutation assessment, and routing assessment types, respectively. Write
$\mathsf{Gen}[X]=\mathtt{Generated}[X]$; $z.\mathtt{value}$ is the declared
projection, not an inferred proof of $X$.
\[
\begin{aligned}
 a &:G\to\mathsf{Gen}[H],&d &:H\to\mathsf{Gen}[C],\\
 r &:C\to\mathsf{Gen}[R],&e &:J\to\mathsf{Gen}[S],\\
 h&=a(g),\quad c=d(h.\mathtt{value}),&v&=r(c.\mathtt{value}),\\
 J&=\{\mathtt{derived\_claim}:C,\ \mathtt{refutation\_assessment}:R\},\\
 s&=e\{\mathtt{derived\_claim}=c.\mathtt{value},\
       \mathtt{refutation\_assessment}=v.\mathtt{value}\}.
\end{aligned}
\]
The relation contracts record assumption, test, and calibration links. In
particular, \emph{Treats\allowbreak Proof\allowbreak Route\allowbreak As\allowbreak Claim\allowbreak Premise} licenses the description
``conditional on this route''; it does not certify that the route holds. The
other two relations attach the challenge and the resulting work assessment to
that commitment. Their source/target records do not supply logical soundness
proofs.

\paragraph{Constructed example.}
To establish a graph-family mixing bound, $a$ proposes a uniform spectral gap.
Under this premise, $d$ derives exponential decay. Testing cycle graphs at $r$
shows their gaps approach zero with size. This refutes the strengthening that
connected graphs have a size-independent gap, not the conditional decay theorem.
The final join supports reframing the
graph family or retargeting the bound, rather than declaring all decay reasoning
false. This is a mathematical illustration, not an observed workflow run.

The native sequence returns a routing assessment with release, reframe,
retarget, or refuted alternatives. The original ERFL additionally switches to
release promotion, conceptual exploration, or target navigation, and stops when
$r$ is refuted. Those transitions, including the early stop before calibration,
are not executable branches in this native sequence.

% SPDX-FileCopyrightText: © 2026 Idiap Research Institute <contact@idiap.ch>
%
% SPDX-FileContributor: Andre Freitas <andre.freitas@idiap.ch>
% Neuro-symbolic AI Group
%
% SPDX-FileContributor: Danilo Gusicuma <danilo.gusicuma@idiap.ch>
% Neuro-symbolic AI Group
%
% SPDX-License-Identifier: MIT
% Generated from step-explanations.json
\paragraph{Step-by-step reconstruction.}
Follow a proposed uniform diffusion bound for du/dt=-Lu over finite connected undirected graphs, using the symmetric combinatorial Laplacian L and mean-zero initial data.
The following is a source-guided illustrative reconstruction. Internal procedures are prescriptions; the native declarations do not establish their execution.
\begin{enumerate}[leftmargin=*,label=\arabic*.]
\item \textbf{Repair the missing premise} (\texttt{route\_\allowbreak{}hypothesis}).
\begin{description}[style=standard,labelindent=0pt,leftmargin=5.5em,labelwidth=5em,labelsep=0.5em,align=left,font=\normalfont\bfseries,itemsep=0.15em,parsep=0pt,topsep=0.3em]
\item[Purpose] Identify a premise sufficient for uniform decay.
\item[Input] Active obligation: decay uniform in graph size
\item[Method] First attempt deduction; locate the absent support, propose small competing premise sets, and attach a falsifier to each.
\item[Check] Compare explanatory coverage and compatibility without treating simplicity as proof.
\item[Output] Route hypothesis: positive common spectral gap
\item[Next] Send this conjectural premise to derivation; incompatible premise sets require revision.
\item[Example] The gap becomes the proposal that every graph in the selected family has first positive eigenvalue at least a common c greater than zero.
\end{description}
\item \textbf{Expose the conditional consequence} (\texttt{derived\_\allowbreak{}claim}).
\begin{description}[style=standard,labelindent=0pt,leftmargin=5.5em,labelwidth=5em,labelsep=0.5em,align=left,font=\normalfont\bfseries,itemsep=0.15em,parsep=0pt,topsep=0.3em]
\item[Purpose] Derive the proposed premise’s conditional payoff.
\item[Input] Proposed common gap; graph-family quantifier
\item[Method] Derive what the added premise would imply while tracking where it is used; do not discharge that premise by using it.
\item[Check] A valid conditional derivation must retain the norm, subspace and common constant.
\item[Output] Exponential-decay claim, conditional on the gap
\item[Next] Refutation receives this claim; a missing spectral lemma remains an obligation.
\item[Example] The proposed gap yields mean-zero norm decay bounded by the initial norm times exp(-ct), conditional on the symmetric Laplacian assumptions.
\end{description}
\item \textbf{Attack the exact strengthening} (\texttt{refutation\_\allowbreak{}assessment}).
\begin{description}[style=standard,labelindent=0pt,leftmargin=5.5em,labelwidth=5em,labelsep=0.5em,align=left,font=\normalfont\bfseries,itemsep=0.15em,parsep=0pt,topsep=0.3em]
\item[Purpose] Test the strengthening without conflating premise and consequence.
\item[Input] Derived decay claim; retained uniform-gap premise
\item[Method] Probe boundary families and seek a counterexample to the implication actually being asserted; distinguish a false premise from an invalid derivation.
\item[Check] A counterexample must satisfy the stated domain and attack the same quantifier.
\item[Output] Refutation finding: connectedness insufficient for a common gap
\item[Next] Pass the targeted finding to calibration; bounded non-refutation supplies no missing proof.
\item[Example] For cycle graphs the first positive eigenvalue is 2-2 cos(2 pi/n), which tends to zero. This defeats a size-independent gap for all connected graphs.
\end{description}
\item \textbf{Choose the next proof task} (\texttt{work\_\allowbreak{}routing}).
\begin{description}[style=standard,labelindent=0pt,leftmargin=5.5em,labelwidth=5em,labelsep=0.5em,align=left,font=\normalfont\bfseries,itemsep=0.15em,parsep=0pt,topsep=0.3em]
\item[Purpose] Choose a sharper task from the localized failure.
\item[Input] Original derived claim + its refutation assessment
\item[Method] Compare the route with the target and identify whether a premise, graph family or bound must change.
\item[Check] Continuation must sharpen the obligation; repeating the same failed premise is not progress.
\item[Output] Routing assessment: bounded reframing or retargeting rationale
\item[Next] Original policy can route onward or stop; the native sequence returns an assessment only.
\item[Example] Retain the conditional decay theorem and replace the unrestricted family claim by the task of finding structural assumptions that imply a uniform gap.
\end{description}
\end{enumerate}
\paragraph{Changed-input contrast.}
If the input instead supplies a proved common spectral-gap bound for a restricted family, the cycle objection no longer defeats that premise. The conditional decay implication can then support work on that family; connectedness alone cannot.


\subsection{Discovery marathon: separating a pattern from its explanation}
\label{app:granular-discovery-marathon-loop}
Discovery here is a bounded comparison of commitments. Scope grounding ($e_0$)
fixes what the run may use. Abduction ($a$) proposes an explanation; induction
($i$) extracts a provisional pattern from that same scope. They are parallel
occurrences because a pattern and an explanatory hypothesis are independently
consumed results. Neither becomes confirmation merely because both were formed.
Model probing ($m$) asks what their combination implies under a model;
contrastive probing ($c$) asks what a discriminating comparison reveals.
Refutation ($r$) combines these different challenges. Calibration ($e_1$)
rejoins the original hypothesis and pattern with the refutation assessment.

\paragraph{What makes a discovery cycle informative.}
The portfolio procedure treats target commitments as fixed constraints and
hypotheses as replaceable. A candidate earns further work by identifying the
commitments it advances and the evidence its next step could change; a novel but
irrelevant branch remains visible without entering the proof queue. The linked
probe procedure asks rivals to predict the same observable under the same
controls before any run. A minimal perturbation of an assumption or carrier then
reveals which dependency the proposed mechanism needs. If rivals cannot be
distinguished, the correct output is non-identifiability, not preference for the
more elaborate model. Negative and boundary outcomes feed branch revision;
repeated failure without new learning forces a change of route. These prescribed
procedures explain the two probe roles, but their current entry gates are richer
than the native hypothesis--pattern join and are not established by its types.

With $X,B,H,P,M,C,R,S$ denoting context, bounded scope, hypothesis, pattern,
model implications, contrast assessment, refutation, and routing:
\[
\begin{aligned}
 e_0 &:X\to\mathsf{Gen}[B],\qquad
 a:B\to\mathsf{Gen}[H],\quad i:B\to\mathsf{Gen}[P],\\
 m &:J_{HP}\to\mathsf{Gen}[M],\qquad
 c:J_{HP}\to\mathsf{Gen}[C],\\
 r &:J_{MC}\to\mathsf{Gen}[R],\qquad
 e_1:J_{HPR}\to\mathsf{Gen}[S],\\
 J_{HP}&=\{\mathtt{discovery\_hypothesis}:H,\ \mathtt{discovery\_pattern}:P\},\\
 J_{MC}&=\{\mathtt{model\_implications}:M,\ \mathtt{contrast\_assessment}:C\},\\
 J_{HPR}&=\{\mathtt{discovery\_hypothesis}:H,\ \mathtt{discovery\_pattern}:P,
          \ \mathtt{refutation\_assessment}:R\}.
\end{aligned}
\]
Both formation calls consume $e_0(x).\mathtt{value}$. Both probes consume the
same named join of $a$'s and $i$'s projected values; $r$ receives the projected
probe outputs; $e_1$ retains the formation outputs alongside $r$'s value.
\emph{Scopes\allowbreak Discovery\allowbreak Formation}, \emph{Grounds\allowbreak Discovery\allowbreak Probe},
\emph{Grounds\allowbreak Discovery\allowbreak Refutation}, and \emph{Calibrates\allowbreak Discovery\allowbreak Routing}
respectively identify these four evidential dependencies, without establishing
that the model or comparison is adequate.

\paragraph{Constructed example.}
Within a fixed sensor dataset, induction finds a drifting offset and abduction
proposes either thermal sensitivity or clock drift. A model predicts offset
changes under temperature variation; the contrast holds temperature fixed while
varying elapsed time. If both hypotheses make the same prediction under available
controls, no run selects between them: the next task is a discriminating control.
If the offset changes only with elapsed time, thermal sensitivity is weakened
within that setup. The portfolio preserves the calibration target and records
which observation could reverse the decision.

The native graph has two parallel stages and returns Promotable, Reorient,
TargetReset, or Stalled. The original routes the first three to release,
conceptual exploration, or navigation, and stops on a stalled reasoning run.
It supplies the cross-cycle control policy; native output alternatives alone do
not execute a marathon or establish termination.

% SPDX-FileCopyrightText: © 2026 Idiap Research Institute <contact@idiap.ch>
%
% SPDX-FileContributor: Andre Freitas <andre.freitas@idiap.ch>
% Neuro-symbolic AI Group
%
% SPDX-FileContributor: Danilo Gusicuma <danilo.gusicuma@idiap.ch>
% Neuro-symbolic AI Group
%
% SPDX-License-Identifier: MIT
% Generated from step-explanations.json
\paragraph{Step-by-step reconstruction.}
A constructed sensor record contains temperature, elapsed time and a drifting offset. Thermal sensitivity and clock drift are rival explanations.
The following is a source-guided illustrative reconstruction. Internal procedures are prescriptions; the native declarations do not establish their execution.
\begin{enumerate}[leftmargin=*,label=\arabic*.]
\item \textbf{Bound the search} (\texttt{search\_\allowbreak{}scope}).
\begin{description}[style=standard,labelindent=0pt,leftmargin=5.5em,labelwidth=5em,labelsep=0.5em,align=left,font=\normalfont\bfseries,itemsep=0.15em,parsep=0pt,topsep=0.3em]
\item[Purpose] Bound admissible evidence and comparisons.
\item[Input] Sensor context; measurement constraints; calibration target
\item[Method] Separate available evidence from desired measurements and delimit the regime in which candidate explanations must operate.
\item[Check] A useful scope retains the target and exposes missing controls.
\item[Output] Bounded sensor record and admissible comparisons
\item[Next] Both formation branches consume this scope; insufficient controls constrain subsequent conclusions.
\item[Example] Freeze the measured offset, temperature and elapsed-time columns; note that no intervention has yet separated thermal and temporal effects.
\end{description}
\item \textbf{Propose a mechanism} (\texttt{discovery\_\allowbreak{}hypothesis}).
\begin{description}[style=standard,labelindent=0pt,leftmargin=5.5em,labelwidth=5em,labelsep=0.5em,align=left,font=\normalfont\bfseries,itemsep=0.15em,parsep=0pt,topsep=0.3em]
\item[Purpose] Propose mechanisms with distinguishable predictions.
\item[Input] Bounded sensor record, shared with pattern induction
\item[Method] Construct rival dependencies and specify the observable on which they could disagree, rather than naming a cause without a test.
\item[Check] A candidate must address the calibration target and admit a falsifier.
\item[Output] Candidate thermal mechanism; clock-drift rival
\item[Next] Both probes receive the hypothesis with the pattern; unsupported mechanisms remain candidates.
\item[Example] Thermal sensitivity predicts offset changes with temperature at fixed elapsed time; clock drift predicts changes with elapsed time at fixed temperature.
\end{description}
\item \textbf{Extract the observed regularity} (\texttt{discovery\_\allowbreak{}pattern}).
\begin{description}[style=standard,labelindent=0pt,leftmargin=5.5em,labelwidth=5em,labelsep=0.5em,align=left,font=\normalfont\bfseries,itemsep=0.15em,parsep=0pt,topsep=0.3em]
\item[Purpose] Describe observed variation without inferring causation.
\item[Input] Bounded observations; no hypothesis output
\item[Method] Compare offsets across the observed conditions and state the trend together with its observational limits.
\item[Check] The pattern must be supported within the record and retain confounding.
\item[Output] Provisional offset trend, with confounding retained
\item[Next] Join this pattern with the hypothesis for both probes; it is not independent confirmation.
\item[Example] The record shows increasing offset while both temperature and time increase. This trend fits both candidate mechanisms and selects neither.
\end{description}
\item \textbf{Derive model predictions} (\texttt{model\_\allowbreak{}implications}).
\begin{description}[style=standard,labelindent=0pt,leftmargin=5.5em,labelwidth=5em,labelsep=0.5em,align=left,font=\normalfont\bfseries,itemsep=0.15em,parsep=0pt,topsep=0.3em]
\item[Purpose] Derive comparable predictions under explicit assumptions.
\item[Input] Discovery hypothesis + provisional observed pattern
\item[Method] Specify variables, controls and measurement assumptions; preregister which outcomes favor, weaken or fail to distinguish each rival, including an inconclusive region, before deriving comparable predictions.
\item[Check] Predictions must concern the same observable and regime as the target.
\item[Output] Conditional thermal and temporal model predictions
\item[Next] Refutation later receives these implications with the contrast assessment; model failure remains informative.
\item[Example] Under thermal sensitivity, holding temperature fixed removes the predicted thermal offset change; under clock drift the offset can still grow with time.
\end{description}
\item \textbf{Design the discriminating contrast} (\texttt{contrast\_\allowbreak{}assessment}).
\begin{description}[style=standard,labelindent=0pt,leftmargin=5.5em,labelwidth=5em,labelsep=0.5em,align=left,font=\normalfont\bfseries,itemsep=0.15em,parsep=0pt,topsep=0.3em]
\item[Purpose] Identify a control that separates the rivals.
\item[Input] Discovery hypothesis + pattern, shared with model probe
\item[Method] Perturb one relevant assumption or control while preserving the others; explain which dependency the contrast tests.
\item[Check] A contrast is discriminating only if rivals predict different outcomes under justified controls.
\item[Output] Contrast assessment: informative control or non-identifiability
\item[Next] Send this assessment to refutation; unavailable controls require a new measurement plan.
\item[Example] Compare readings at fixed temperature but different elapsed times. Continued drift challenges a purely thermal account if other conditions remain stable.
\end{description}
\item \textbf{Combine the challenges} (\texttt{refutation\_\allowbreak{}assessment}).
\begin{description}[style=standard,labelindent=0pt,leftmargin=5.5em,labelwidth=5em,labelsep=0.5em,align=left,font=\normalfont\bfseries,itemsep=0.15em,parsep=0pt,topsep=0.3em]
\item[Purpose] Locate what survives both challenges.
\item[Input] Model implications + contrast assessment; scoped assumptions
\item[Method] Compare predicted differences with the available contrast and localize inconsistency; do not count two analyses of one record as two independent experiments.
\item[Check] A negative verdict needs a particular violated prediction, not an unfavorable label.
\item[Output] Refutation assessment: proposed falsifier, no observed refutation
\item[Next] Calibration receives this finding with the original hypothesis and pattern.
\item[Example] If a future fixed-temperature test finds drift, it would defeat the thermal-only prediction under its control assumptions. The current confounded record supplies no such finding.
\end{description}
\item \textbf{Allocate the next discovery action} (\texttt{routing\_\allowbreak{}assessment}).
\begin{description}[style=standard,labelindent=0pt,leftmargin=5.5em,labelwidth=5em,labelsep=0.5em,align=left,font=\normalfont\bfseries,itemsep=0.15em,parsep=0pt,topsep=0.3em]
\item[Purpose] Choose continuation, reorientation or reset.
\item[Input] Hypothesis + pattern + refutation; retained calibration target
\item[Method] Compare target relevance and what the next action could change; preserve alternatives that remain indistinguishable.
\item[Check] Continuation needs new learning or a sharper test, not repeated failure.
\item[Output] Routing assessment: remaining mechanism question and justification
\item[Next] Original cross-cycle routing is a policy; the native return does not execute a marathon.
\item[Example] Retain the trend and both rival explanations; request the fixed-temperature measurement. Demote thermal-only only if the predicted falsifier is subsequently observed under the stated controls.
\end{description}
\end{enumerate}
\paragraph{Changed-input contrast.}
If temperature and elapsed time are perfectly confounded, neither probe identifies the cause. The proper next action is a discriminating measurement, not promotion of the more elaborate explanation.


\subsection{Conceptual repair: turning an obstruction into a testable bridge}
\label{app:granular-conceptual-hypothesis-exploration-loop}
This workflow changes the representation used to attack a problem. Obstruction
assessment ($e_0$) identifies the structural pressure behind a failure. Mediation
formation ($a$) proposes what kind of concept could resolve it. Analogical
transfer ($t$) builds a correspondence rather than merely naming a similar
subject. Deduction ($d$) exposes a consequence of that bridge. Refutation ($r$)
tests the consequence, and routing ($e_1$) evaluates the bridge in light of that
test. The distinction between mediation and bridge prevents a promising analogy
from being treated as an already specified transport of reasoning.

\paragraph{What is revised, and why.}
The linked bridge-forging procedure searches for a missing carrier, invariant,
representation, interface, or transfer law where two structures nearly meet.
It must state why existing objects fail, how the new object would be constructed,
which invariants survive, and its first failure regime. Defining the desired
conclusion into the bridge is explicitly disallowed. The route ranker then
compares target contribution, validation strength, risk, and tractability; a
beautiful bridge without proof-form obligations stays outside the proof queue.
Its feedback decision distinguishes informative adaptation from backtracking,
pivoting, and quarantine, and requires the selected machinery to control the
actual bottleneck. The workflow contract retains failure cartography, mediation
candidates, and validation packets for this decision. This describes prescribed
conceptual work around the native ports, not a recovered implementation of their
string-bearing values.

Let $F,P,H,B,C,R,S$ name failure state, obstruction pressure, mediation
hypothesis, bridge, consequence, refutation, and refinement status.
\[
\begin{aligned}
 e_0 &:F\to\mathsf{Gen}[P],&a &:P\to\mathsf{Gen}[H],\\
 t &:H\to\mathsf{Gen}[B],&d &:B\to\mathsf{Gen}[C],\\
 r &:C\to\mathsf{Gen}[R],&e_1 &:J_{BR}\to\mathsf{Gen}[S],\\
 p&=e_0(f),\ h=a(p.\mathtt{value}),\ b=t(h.\mathtt{value}),\\
 q&=d(b.\mathtt{value}),\ v=r(q.\mathtt{value}),\\
 J_{BR}&=\{\mathtt{conceptual\_bridge}:B,\ \mathtt{refutation\_assessment}:R\},\\
 s&=e_1\bigl\{\mathtt{conceptual\_bridge}=b.\mathtt{value},\\
  &\hspace{3em}\mathtt{refutation\_assessment}=v.\mathtt{value}\bigr\}.
\end{aligned}
\]
The relation sequence bounds obstruction pressure, transfers mediation into a
bridge, treats the bridge as a consequence premise, tests that consequence,
and calibrates routing. In particular, the transfer relation records the proposed
correspondence; the assumption relation keeps its consequences conditional.
Neither declaration asserts that analogy preserves all relevant structure.

\paragraph{Constructed example.}
A convergence argument for a physical state fails because a norm counts changes
of gauge as changes of observable state. Obstruction assessment identifies that
identity mismatch. Mediation proposes a quotient by gauge equivalence; the bridge
must construct the projection and show that dynamics descend to the quotient.
Deduction can then formulate convergence of equivalence classes. If two equivalent
representatives evolve into inequivalent classes, refutation exposes failure of
well-defined dynamics. The next task is to repair that compatibility law, not to
rename the quotient as a solution. This constructed case illustrates how a
representation change generates a new, exact proof burden.

Native output distinguishes WorkReady, ReadyForReleasePromotion, and Refuted.
Original ERFL sends work-ready repairs back to abductive proof work, release
candidates to promotion, and stops on refutation. The native straight sequence
contains neither these switches nor the original early-stop guard. Its final
join retains the bridge and assessment; it has no separately typed retained port
for the derived consequence or original failure state.

% SPDX-FileCopyrightText: © 2026 Idiap Research Institute <contact@idiap.ch>
%
% SPDX-FileContributor: Andre Freitas <andre.freitas@idiap.ch>
% Neuro-symbolic AI Group
%
% SPDX-FileContributor: Danilo Gusicuma <danilo.gusicuma@idiap.ch>
% Neuro-symbolic AI Group
%
% SPDX-License-Identifier: MIT
% Generated from step-explanations.json
\paragraph{Step-by-step reconstruction.}
A state-convergence argument uses a norm that distinguishes gauge representatives even when they represent the same observable state.
The following is a source-guided illustrative reconstruction. Internal procedures are prescriptions; the native declarations do not establish their execution.
\begin{enumerate}[leftmargin=*,label=\arabic*.]
\item \textbf{Locate the representational obstruction} (\texttt{obstruction\_\allowbreak{}pressure}).
\begin{description}[style=standard,labelindent=0pt,leftmargin=5.5em,labelwidth=5em,labelsep=0.5em,align=left,font=\normalfont\bfseries,itemsep=0.15em,parsep=0pt,topsep=0.3em]
\item[Purpose] Distinguish conceptual obstruction from unfinished calculation.
\item[Input] Failed argument; norm’s identity distinction
\item[Method] Trace the failed implication to the role played by the carrier, invariant or equivalence relation.
\item[Check] An obstruction diagnosis must explain why existing objects fail at the named boundary.
\item[Output] Obstruction diagnosis: representative-dependent convergence criterion
\item[Next] Mediation consumes this diagnosis; a merely missing calculation should remain a derivation task.
\item[Example] Two gauge-equivalent representatives can be far apart in the chosen norm, so representative convergence need not express physical-state convergence.
\end{description}
\item \textbf{Propose the mediating object} (\texttt{mediation\_\allowbreak{}hypothesis}).
\begin{description}[style=standard,labelindent=0pt,leftmargin=5.5em,labelwidth=5em,labelsep=0.5em,align=left,font=\normalfont\bfseries,itemsep=0.15em,parsep=0pt,topsep=0.3em]
\item[Purpose] Propose a carrier for the intended question.
\item[Input] Representative-dependent identity and convergence obstruction
\item[Method] Propose a quotient or other mediator, naming what it would preserve and its first possible failure regime.
\item[Check] The desired convergence theorem cannot be hidden inside the mediator's definition.
\item[Output] Mediation hypothesis: gauge-equivalence classes
\item[Next] Transfer must turn this proposal into specified maps; a name alone is insufficient.
\item[Example] Replace the proposed carrier of raw representatives by classes under gauge equivalence, while leaving existence of suitable dynamics as an obligation.
\end{description}
\item \textbf{Specify the bridge} (\texttt{conceptual\_\allowbreak{}bridge}).
\begin{description}[style=standard,labelindent=0pt,leftmargin=5.5em,labelwidth=5em,labelsep=0.5em,align=left,font=\normalfont\bfseries,itemsep=0.15em,parsep=0pt,topsep=0.3em]
\item[Purpose] Specify the mediator’s maps and compatibility burden.
\item[Input] Quotient proposal; intended observable interpretation
\item[Method] Specify projection from representatives to classes and the compatibility needed for dynamics to descend.
\item[Check] Equivalent representatives must induce the same class evolution for the bridge to be well defined.
\item[Output] Projection bridge; open dynamics-compatibility obligation
\item[Next] Deduction uses this bridge conditionally; an absent compatibility law blocks unconditional transport.
\item[Example] For representative evolution F, require the class of F(x) to equal the class of F(y) whenever x and y are gauge equivalent.
\end{description}
\item \textbf{Derive a bridge consequence} (\texttt{bridge\_\allowbreak{}consequence}).
\begin{description}[style=standard,labelindent=0pt,leftmargin=5.5em,labelwidth=5em,labelsep=0.5em,align=left,font=\normalfont\bfseries,itemsep=0.15em,parsep=0pt,topsep=0.3em]
\item[Purpose] State the bridge’s conditional consequence.
\item[Input] Mediating bridge; open compatibility assumptions
\item[Method] Formulate the transported statement on the new carrier and retain each bridge premise.
\item[Check] The consequence must concern equivalence classes without silently restoring representative identity.
\item[Output] Class-level evolution or convergence claim, conditional on its premises
\item[Next] Refutation tests this consequence; convergence estimates remain separate from well-definedness.
\item[Example] A quotient trajectory can be defined from F only if changing the representative leaves its next class unchanged; convergence would then be a class-level question.
\end{description}
\item \textbf{Test preservation} (\texttt{refutation\_\allowbreak{}assessment}).
\begin{description}[style=standard,labelindent=0pt,leftmargin=5.5em,labelwidth=5em,labelsep=0.5em,align=left,font=\normalfont\bfseries,itemsep=0.15em,parsep=0pt,topsep=0.3em]
\item[Purpose] Test whether the representation is well defined.
\item[Input] Bridge consequence; compatibility premise
\item[Method] Search for equivalent inputs whose evolved outputs cease to be equivalent; localize the precise preservation law at fault.
\item[Check] A bridge fails this test when one physical input class receives two distinct output classes.
\item[Output] Refutation assessment: well-definedness defect or bounded surviving claim
\item[Next] Pass the finding to routing with the retained bridge; no found defect is not a proof.
\item[Example] Choose equivalent x and y with inequivalent F(x) and F(y). The proposed class dynamics is ambiguous, independently of any convergence estimate.
\end{description}
\item \textbf{Route the conceptual repair} (\texttt{refinement\_\allowbreak{}routing}).
\begin{description}[style=standard,labelindent=0pt,leftmargin=5.5em,labelwidth=5em,labelsep=0.5em,align=left,font=\normalfont\bfseries,itemsep=0.15em,parsep=0pt,topsep=0.3em]
\item[Purpose] Choose proof work or bridge reconstruction.
\item[Input] Conceptual bridge + its refutation assessment
\item[Method] Compare target contribution, validation strength and the exact preservation defect; distinguish an informative repair from a repeated blocked route.
\item[Check] A work-ready bridge needs explicit obligations that address the bottleneck.
\item[Output] Refinement assessment: compatibility-repair obligation
\item[Next] Original policy routes work or release and stops refutation; native routing is a returned value.
\item[Example] Keep quotienting as a candidate representation but require a compatible evolution law; do not advertise the quotient name as a solved convergence problem.
\end{description}
\end{enumerate}
\paragraph{Changed-input contrast.}
If equivalent representatives evolve into equivalent classes, the proposed quotient passes this well-definedness test. Convergence still requires its own argument; compatibility alone does not establish it.


\subsection{Frontier deepening: selecting before formalizing}
\label{app:granular-conceptual-frontier-deepening}
A concept should become a formal claim only after its discriminating content has
survived an explicit selection step. Frontier grounding ($e_0$) fixes a bounded
search from frozen context. Hypothesis formation ($a$) proposes a candidate
within it. Refutation ($r$) challenges that candidate before selection ($e_1$)
compares its promise with the challenge. Bridge revision ($t$) turns the selected
content into a correspondence usable for formal work; deduction ($d$) expresses
a formal consequence. Readiness calibration ($e_2$) decides whether that claim
is ready for theorem work, not whether it is proved.

\paragraph{The conceptual work hidden by a single hypothesis port.}
The linked frontier procedure opens invention before assigning family labels:
authors work without sibling candidates, then a separate assessor reconstructs
families from primitive changes, mechanism, intervention, prediction, and first
falsifier. Semantic paraphrases collapse; incompatible mechanisms remain live.
Reflection changes identity conditions, exposes internal contradictions and
inherited commitments, separates necessary from regime-relative structure, and
seeks unification that preserves decisive differences. Blind judgment compares
complete candidates on non-compensatory dimensions such as relevance,
protected-success preservation, explanatory depth, correction, and falsifiability;
a plural non-dominated frontier is an admissible outcome. Category or HoTT
representations earn a role only by exposing an actual mediator or transport
obstruction. This prescribed process is much richer than the linear native
projection. There is also a source difference: the skill makes formal routes
conditional, while the named contract lists central-route derivation and formal
claim production as required steps. Neither should be silently read as proof
that every candidate becomes formal-ready.

The local aliases are $X$ (frozen context), $F$ (frontier), $H$ (hypothesis),
$R$ (refutation), $S$ (selection), $B$ (revised bridge), $C$ (formal claim), and
$Q$ (readiness).
\[
\begin{aligned}
 e_0 &:X\to\mathsf{Gen}[F],&a &:F\to\mathsf{Gen}[H],\\
 r &:H\to\mathsf{Gen}[R],&e_1 &:J_{HR}\to\mathsf{Gen}[S],\\
 t &:S\to\mathsf{Gen}[B],&d &:B\to\mathsf{Gen}[C],\\
 e_2 &:C\to\mathsf{Gen}[Q],\\
 J_{HR}&=\{\mathtt{conceptual\_hypothesis}:H,
           \ \mathtt{refutation\_assessment}:R\}.
\end{aligned}
\]
Every sequential call uses the predecessor's $\mathtt{.value}$; selection instead
joins the retained hypothesis value with its refutation value. The search-bound
relation attaches formation to the grounded frontier. The test relation attaches
criticism to the candidate, and the calibration relation is used twice: first
for selection, then for readiness. Bridge revision and premise relations mark
the remaining changes of commitment. Reusing a relation name does not equate
these two calibration judgments.

\paragraph{Constructed example.}
For two effective descriptions with different state variables but matching
measured predictions in a bounded regime, invention may propose a common latent
carrier or a direct translation between observables. Reflection asks which
identity distinctions each proposal erases. A candidate transport must specify
maps, the regime in which they apply, and a test where predictions diverge.
Agreement on existing observations does not establish equivalence, and a latent
carrier without a discriminating consequence receives no automatic preference.
Both candidates can remain on the frontier until a correction or new observable
distinguishes them. The formalization task, when justified, is the specific
transport or compatibility obligation.

The native output offers FormalReady, NotFormalReadyOrNonDominated, or Refuted,
but the native sequence has no switch. Original ERFL routes readiness to theorem
proof, returns a non-ready or non-dominated candidate to exploration, and stops
at a refuted $r$. Moreover, selection, bridge, and claim are nominal string-bearing
types: their signatures do not enforce preservation of the original evidence or
hypothesis identity. Such stronger indexed obligations must be proposed
explicitly rather than attributed to this native declaration.

% SPDX-FileCopyrightText: © 2026 Idiap Research Institute <contact@idiap.ch>
%
% SPDX-FileContributor: Andre Freitas <andre.freitas@idiap.ch>
% Neuro-symbolic AI Group
%
% SPDX-FileContributor: Danilo Gusicuma <danilo.gusicuma@idiap.ch>
% Neuro-symbolic AI Group
%
% SPDX-License-Identifier: MIT
% Generated from step-explanations.json
\paragraph{Step-by-step reconstruction.}
Two effective descriptions match existing measurements in a bounded regime but use different state variables. A latent carrier and a direct observable translation are candidate bridges.
The following is a source-guided illustrative reconstruction. Internal procedures are prescriptions; the native declarations do not establish their execution.
\begin{enumerate}[leftmargin=*,label=\arabic*.]
\item \textbf{Freeze the conceptual pressure} (\texttt{conceptual\_\allowbreak{}frontier}).
\begin{description}[style=standard,labelindent=0pt,leftmargin=5.5em,labelwidth=5em,labelsep=0.5em,align=left,font=\normalfont\bfseries,itemsep=0.15em,parsep=0pt,topsep=0.3em]
\item[Purpose] Fix protected successes and the unresolved distinction.
\item[Input] Frozen descriptions; bounded empirical agreement
\item[Method] Separate target commitments from replaceable concepts and identify what a new bridge must explain.
\item[Check] The frontier must retain regime limits and the successes a replacement may not discard.
\item[Output] Bounded frontier: relation between descriptions; protected distinctions
\item[Next] Hypothesis formation consumes this frontier; ordinary missing calculation does not justify a new ontology.
\item[Example] Record agreement on current observables, different state variables, and an open question about predictions beyond that observable set.
\end{description}
\item \textbf{Invent alternatives} (\texttt{conceptual\_\allowbreak{}hypothesis}).
\begin{description}[style=standard,labelindent=0pt,leftmargin=5.5em,labelwidth=5em,labelsep=0.5em,align=left,font=\normalfont\bfseries,itemsep=0.15em,parsep=0pt,topsep=0.3em]
\item[Purpose] Generate meaningfully different explanatory structures.
\item[Input] Bounded frontier; protected commitments
\item[Method] The linked skill prescribes 8--12 isolated raw candidates, blinded reconstruction of at least six semantic families unless verifier-replayed constraint evidence justifies fewer, and six reflective transformations. Propose primitive, mediator or transport changes and candidate-specific falsifiers; this two-candidate example illustrates one distinction only.
\item[Check] Semantic paraphrases do not create new families; each alternative must change a meaningful dependency.
\item[Output] Candidate explanations: latent mediation or direct observable translation
\item[Next] Refutation examines candidate content; a nominal hypothesis port does not implement the larger portfolio procedure.
\item[Example] A latent carrier attempts to explain both state descriptions, while direct translation relates selected observables without asserting a common underlying state space.
\end{description}
\item \textbf{Challenge a candidate's distinctions} (\texttt{refutation\_\allowbreak{}assessment}).
\begin{description}[style=standard,labelindent=0pt,leftmargin=5.5em,labelwidth=5em,labelsep=0.5em,align=left,font=\normalfont\bfseries,itemsep=0.15em,parsep=0pt,topsep=0.3em]
\item[Purpose] Test what the proposed bridge fails to preserve.
\item[Input] Candidate explanation; claimed regime
\item[Method] Identify a decisive observable, intervention or identity distinction and ask whether the proposal preserves it.
\item[Check] A criticism must change a dependency or reveal an unmet prediction, not merely request different wording.
\item[Output] Refutation finding: unresolved discrimination or specific candidate defect
\item[Next] Selection receives the original hypothesis with this finding, preserving its identity.
\item[Example] Ask whether two states merged by the latent carrier could disagree on a newly measured observable. If so, that merging needs justification or repair.
\end{description}
\item \textbf{Select without forcing consensus} (\texttt{hypothesis\_\allowbreak{}selection}).
\begin{description}[style=standard,labelindent=0pt,leftmargin=5.5em,labelwidth=5em,labelsep=0.5em,align=left,font=\normalfont\bfseries,itemsep=0.15em,parsep=0pt,topsep=0.3em]
\item[Purpose] Select defensible content without forcing preference.
\item[Input] Hypothesis + its own refutation assessment
\item[Method] The linked skill prescribes preregistered non-compensatory dimensions, at least three separately configured blinded judges and sham/rival controls. Compare complete candidates and retain a nondominated set; missing judges, warrants or failed controls yield a typed null rather than a forced preference.
\item[Check] Failure on a required dimension cannot be offset by elegance elsewhere.
\item[Output] Selected content; unresolved alternatives where warranted
\item[Next] Bridge revision consumes selected content; non-ready plural results have an original exploration route.
\item[Example] Keep both latent mediation and direct translation if current evidence cannot distinguish them; make the missing observable test part of their retained burden.
\end{description}
\item \textbf{Make a selected bridge explicit} (\texttt{conceptual\_\allowbreak{}bridge}).
\begin{description}[style=standard,labelindent=0pt,leftmargin=5.5em,labelwidth=5em,labelsep=0.5em,align=left,font=\normalfont\bfseries,itemsep=0.15em,parsep=0pt,topsep=0.3em]
\item[Purpose] Specify a useful formal correspondence.
\item[Input] Selected content; open identity and regime obligations
\item[Method] Specify maps, carriers and preservation conditions only where formal representation adds a concrete obligation.
\item[Check] Category or type language must expose a mediator or transport obstruction, not supply prestige.
\item[Output] Revised bridge: specified compatibility obligation
\item[Next] Deduction uses this bridge; linked procedures make formal routes conditional although the native sequence is linear.
\item[Example] For the direct-translation candidate, specify an observable map and demand agreement of predictions over the chosen regime, without assuming full state-space equivalence.
\end{description}
\item \textbf{State the exact formal burden} (\texttt{formal\_\allowbreak{}claim}).
\begin{description}[style=standard,labelindent=0pt,leftmargin=5.5em,labelwidth=5em,labelsep=0.5em,align=left,font=\normalfont\bfseries,itemsep=0.15em,parsep=0pt,topsep=0.3em]
\item[Purpose] Formulate the bridge’s precise proof obligation.
\item[Input] Revised map; preservation conditions
\item[Method] Express the compatibility statement with explicit domain and quantifiers, retaining assumptions as premises.
\item[Check] The claim must not strengthen regime-limited prediction agreement to unrestricted equivalence.
\item[Output] Formal map-compatibility claim; retained premises and regime
\item[Next] Readiness assesses this statement; the act of stating it does not discharge it.
\item[Example] The obligation becomes prediction preservation for the specified observable map on the bounded regime; unseen observables remain outside the claim.
\end{description}
\item \textbf{Assess readiness for theorem work} (\texttt{calibrated\_\allowbreak{}candidate}).
\begin{description}[style=standard,labelindent=0pt,leftmargin=5.5em,labelwidth=5em,labelsep=0.5em,align=left,font=\normalfont\bfseries,itemsep=0.15em,parsep=0pt,topsep=0.3em]
\item[Purpose] Assess proof readiness while retaining unresolved choices.
\item[Input] Formal claim; domain; open obligations
\item[Method] Check that objects and the required implication are definite; distinguish readiness from an established theorem.
\item[Check] An unresolved choice of carrier or target statement prevents an unconditional readiness judgment.
\item[Output] Calibrated candidate: ready, non-ready or unresolved plural status
\item[Next] Original policy can send ready claims to theorem work; the native sequence supplies no switch.
\item[Example] If the observable map and its regime are specified, queue its compatibility lemma; if the carrier remains ambiguous, return to conceptual selection.
\end{description}
\end{enumerate}
\paragraph{Changed-input contrast.}
If a new observable distinguishes the candidates, selection must incorporate that difference. Agreement on the original measurements alone warrants neither equivalence nor a unique winner.


\subsection{Evidence-to-proof onramp: finding the obligation worth proving}
\label{app:granular-evidence-to-proof-route-onramp}
The onramp separates proposing a useful route from selecting an exact entrypoint.
Grounding ($e_0$) anchors a source bundle in evidence. Route abduction ($a_0$)
proposes how that evidence could support proof work. Deduction ($d_0$) makes a
consequence inspectable, and calibration ($e_1$) assesses it. A second abductive
occurrence ($a_1$) then revises the proof-obligation entrypoint: assessment may
show that a different lemma, rather than the original ambitious target, is the
tractable next step. Refutation ($r_0$) challenges that entrypoint; final
calibration ($e_2$) preserves it with its falsification record.

\paragraph{How an entrypoint is found.}
The named contract orders ingestion and theorem normalization before conjecture
grounding and connection search. The connector distinguishes established
statement dependencies from proposed cross-component links, using symbol roles,
method overlap, and candidate reductions to expose missing bridges. The hub audit
asks whether a highly connected definition is being mistaken for mathematical
progress. Pathfinding then anchors each agenda on one statement, attaches its
supporting hypotheses and links, and gives it a stop or pruning condition. A
regime-crossing route additionally needs an audit of the transfer; the final
handoff retains the route together with its falsification evidence. Consequently,
route selection concerns the missing implication that would advance the target,
not simply the most connected node. These are ordered contract and prescribed
skill procedures; the seven native occurrences are their coarse conceptual
projection, not an execution trace of all those steps.

For source $X$, ground $G$, route $H$, consequence $C$, assessment $A$, entrypoint
$P$, falsification $R$, and readiness $S$:
\[
\begin{aligned}
 e_0 &:X\to\mathsf{Gen}[G],&a_0 &:G\to\mathsf{Gen}[H],\\
 d_0 &:H\to\mathsf{Gen}[C],&e_1 &:C\to\mathsf{Gen}[A],\\
 a_1 &:A\to\mathsf{Gen}[P],&r_0 &:P\to\mathsf{Gen}[R],\\
 e_2 &:J_{PR}\to\mathsf{Gen}[S],\\
 J_{PR}&=\{\mathtt{proof\_entrypoint}:P,
           \ \mathtt{falsification\_assessment}:R\}.
\end{aligned}
\]
After grounding, successive unary calls pass explicit $\mathtt{.value}$
projections along the sequence; $e_2$ joins the projected output of $a_1$ with
the projected output of $r_0$ in the named fields.
Relations identify grounding, conditional consequence, calibration,
entrypoint revision, and testing; calibration is reused at both assessments.
The readiness record contains a disposition, expert dossier, falsification
assessment, and hub-bias audit. The audit is an output field, not an additional
declared reasoning occurrence.

\paragraph{Constructed example.}
A knowledge base connects a circuit lower-bound target to a highly central
complexity-class definition and to a less connected lemma translating proofs into
circuits. The hub audit prevents the definition's degree from deciding the route.
Pathfinding instead anchors on the transfer lemma and asks whether it preserves
uniformity. A counterexample search tests whether the translation relies on
nonuniform advice; if it does, that route cannot support the uniform target
without a new lemma. The dossier therefore carries an explicit uniformity
obligation or a stopped route, rather than a graph-centrality score presented as
proof progress. This is a constructed route diagnosis, not a lower-bound result.

Original ERFL switches a ready route to theorem delivery, sends conceptual gaps
to exploration, and stops on refutation. Native comments describe these
handoffs, including an adapter to \texttt{ProofQueueItem}, as non-executable.
The native sequence itself returns the assessment; it neither implements that
adapter nor enforces the early stop. The nominal ID-bearing intermediate records
also leave evidence-preservation obligations outside their signatures.

% SPDX-FileCopyrightText: © 2026 Idiap Research Institute <contact@idiap.ch>
%
% SPDX-FileContributor: Andre Freitas <andre.freitas@idiap.ch>
% Neuro-symbolic AI Group
%
% SPDX-FileContributor: Danilo Gusicuma <danilo.gusicuma@idiap.ch>
% Neuro-symbolic AI Group
%
% SPDX-License-Identifier: MIT
% Generated from step-explanations.json
\paragraph{Step-by-step reconstruction.}
A lower-bound knowledge base contains a highly connected complexity-class definition and a less connected proposed proof-to-circuit translation. The target requires uniformity.
The following is a source-guided illustrative reconstruction. Internal procedures are prescriptions; the native declarations do not establish their execution.
\begin{enumerate}[leftmargin=*,label=\arabic*.]
\item \textbf{Ground the route evidence} (\texttt{e0}).
\begin{description}[style=standard,labelindent=0pt,leftmargin=5.5em,labelwidth=5em,labelsep=0.5em,align=left,font=\normalfont\bfseries,itemsep=0.15em,parsep=0pt,topsep=0.3em]
\item[Purpose] Ground statements before proposing proof connections.
\item[Input] Source bundle; target; source references
\item[Method] Normalize theorem bodies and retain hypotheses, provenance and the status of each imported claim.
\item[Check] Only licensed statements may serve as established premises; topic proximity is insufficient.
\item[Output] Grounded definitions; proposed translation claims with status
\item[Next] Route formation consumes this evidence; ungrounded statements remain conjectural or excluded.
\item[Example] Record the uniform lower-bound target, the central definition, and a translation candidate whose uniformity condition has not been checked.
\end{description}
\item \textbf{Propose a useful connection} (\texttt{a0}).
\begin{description}[style=standard,labelindent=0pt,leftmargin=5.5em,labelwidth=5em,labelsep=0.5em,align=left,font=\normalfont\bfseries,itemsep=0.15em,parsep=0pt,topsep=0.3em]
\item[Purpose] Identify a target-relevant missing implication.
\item[Input] Grounded statements; dependency context
\item[Method] Compare candidate reductions and cross-component links using symbol roles and method overlap, separating proposed edges from established dependencies.
\item[Check] A route must supply proof relevance beyond graph centrality.
\item[Output] Route hypothesis: candidate proof-to-circuit translation
\item[Next] Deduction exposes its consequence; the selected bridge is still hypothetical.
\item[Example] Choose the less central translation because, if uniform, it could connect a proof lower bound to the circuit statement required by the target.
\end{description}
\item \textbf{Expose the route's consequence} (\texttt{d0}).
\begin{description}[style=standard,labelindent=0pt,leftmargin=5.5em,labelwidth=5em,labelsep=0.5em,align=left,font=\normalfont\bfseries,itemsep=0.15em,parsep=0pt,topsep=0.3em]
\item[Purpose] Expose the candidate bridge’s mathematical payoff.
\item[Input] Translation hypothesis; retained conditions
\item[Method] Derive the proposed implication while tracking the target quantifier that the translation must preserve.
\item[Check] Do not replace a uniform target with a nonuniform consequence.
\item[Output] Translation consequence, conditional on required uniformity
\item[Next] Assessment checks this consequence's relevance and missing support.
\item[Example] The route would transfer the chosen lower-bound obligation only if one uniform construction supplies the translated objects across input sizes.
\end{description}
\item \textbf{Assess the consequence} (\texttt{e1}).
\begin{description}[style=standard,labelindent=0pt,leftmargin=5.5em,labelwidth=5em,labelsep=0.5em,align=left,font=\normalfont\bfseries,itemsep=0.15em,parsep=0pt,topsep=0.3em]
\item[Purpose] Identify the decisive gap in the payoff.
\item[Input] Derived consequence; uniformity assumption
\item[Method] Compare the consequence to the original target and inspect whether graph prominence concealed a missing theorem.
\item[Check] A highly connected definition does not compensate for an unproved translation property.
\item[Output] Consequence assessment: missing uniformity-preservation support
\item[Next] Entrypoint revision consumes this diagnosis instead of repeating the full ambitious target.
\item[Example] The definition contributes vocabulary, but the route's progress depends on proving that translation does not use size-specific advice.
\end{description}
\item \textbf{Revise the entrypoint} (\texttt{a1}).
\begin{description}[style=standard,labelindent=0pt,leftmargin=5.5em,labelwidth=5em,labelsep=0.5em,align=left,font=\normalfont\bfseries,itemsep=0.15em,parsep=0pt,topsep=0.3em]
\item[Purpose] Reduce the route to a testable obligation.
\item[Input] Consequence assessment; localized uniformity gap
\item[Method] Anchor an agenda on one statement, attach its hypotheses and links, and specify a stop or pruning condition.
\item[Check] The entrypoint must directly control the missing implication.
\item[Output] Entrypoint hypothesis: uniformity-preservation lemma
\item[Next] Refutation challenges this entrypoint before a proof handoff is justified.
\item[Example] Replace the full lower-bound demand with the lemma that the proposed translation uses no nonuniform advice as input size varies.
\end{description}
\item \textbf{Challenge the entrypoint} (\texttt{r0}).
\begin{description}[style=standard,labelindent=0pt,leftmargin=5.5em,labelwidth=5em,labelsep=0.5em,align=left,font=\normalfont\bfseries,itemsep=0.15em,parsep=0pt,topsep=0.3em]
\item[Purpose] Test the selected lemma in its stated regime.
\item[Input] Uniformity-preservation entrypoint; hypotheses
\item[Method] Search for a translation instance requiring advice or a regime change; distinguish an actual counterexample from a missing argument.
\item[Check] A falsifier must attack the selected lemma, not a different complexity claim.
\item[Output] Falsification finding: advice dependence or unresolved uniformity proof
\item[Next] Final readiness joins this finding with the same entrypoint.
\item[Example] A translation that chooses unrelated advice at each input size defeats the proposed uniformity lemma, leaving any nonuniform statement separate.
\end{description}
\item \textbf{Bound the proof handoff} (\texttt{e2}).
\begin{description}[style=standard,labelindent=0pt,leftmargin=5.5em,labelwidth=5em,labelsep=0.5em,align=left,font=\normalfont\bfseries,itemsep=0.15em,parsep=0pt,topsep=0.3em]
\item[Purpose] Assess a proof task with its criticism intact.
\item[Input] Entrypoint hypothesis + its falsification assessment
\item[Method] Retain the route dossier, hub-bias diagnosis and transfer burden when judging readiness.
\item[Check] Readiness cannot arise solely from connectivity or a favorable label.
\item[Output] Readiness assessment: bounded proof task or reorientation rationale
\item[Next] Original routing needs a separate ProofQueueItem adapter; the native return does not construct it.
\item[Example] Hand off a precise unresolved uniformity obligation if viable, or stop the advice-dependent route without claiming a lower bound.
\end{description}
\end{enumerate}
\paragraph{Changed-input contrast.}
If the translation is independently shown to preserve uniformity, it can support a suitable proof entrypoint. Merely increasing the centrality of the definition never supplies that implication.


\subsection{Theorem delivery: preserving the exact claim through challenge}
\label{app:granular-theorem-proof-delivery}
The theorem workflow begins after a proof obligation has been chosen. Deduction
($d_0$) constructs an exact theorem claim from explicit premises. Refutation
($r_0$) tests that claim rather than a weaker paraphrase. Calibration ($e_0$)
combines the retained claim and refutation assessment into a disposition with a
formal-proof gate and, where applicable, a failure certificate. This is a
three-occurrence closure attempt: choosing a new hypothesis is deliberately not
a local occurrence. A failed attempt can instead supply conceptual work with an
exact statement and its unresolved obligations.

\paragraph{What happens inside a proof attempt.}
The linked builder orders lemmas without circular dependencies and records where
each premise enters. A parameter ledger propagates constants, exponents,
quantifiers, norms, and validity regimes so that a locally plausible step cannot
silently change the theorem. When a lemma blocks, the conditional gap-splitting
procedure identifies whether the cause is a missing estimate, transfer, wrong
quantifier, model mismatch, or false statement. It actively searches small,
degenerate, proxy, and adversarial cases before classifying the problem as missing
technique, then permits one explicit repair before choosing a new disposition.
The gate distinguishes an unconditional proof from an explicitly conditional
theorem and rejects silent gap closure. The contract therefore retains obligations
even when exit is allowed. These prescribed proof procedures supply substance to
``derive, test, calibrate'' without asserting that the three native operations
execute them or that a status value constitutes a formal proof certificate.

Write $P,C,R,S$ for explicit premises, exact claim, refutation, and disposition.
The input $q$ has native type \texttt{ProofQueueItem}.
\[
\begin{aligned}
 d_0 &:P\to\mathsf{Gen}[C],\qquad
 r_0:C\to\mathsf{Gen}[R],\qquad e_0:J_{CR}\to\mathsf{Gen}[S],\\
 c&=d_0(q.\mathtt{premises}),\qquad v=r_0(c.\mathtt{value}),\\
 J_{CR}&=\{\mathtt{theorem\_claim}:C,
           \ \mathtt{refutation\_assessment}:R\},\\
 s&=e_0\{\mathtt{theorem\_claim}=c.\mathtt{value},\
         \mathtt{refutation\_assessment}=v.\mathtt{value}\}.
\end{aligned}
\]
\emph{Tests\allowbreak Exact\allowbreak Theorem\allowbreak Claim} connects the challenge to the generated claim;
\emph{Calibrates\allowbreak Theorem\allowbreak Proof\allowbreak Disposition} connects the joint evidence to status.
The claim record explicitly retains premises and formal obligations. The output
retains claim, refutation, gate, and an optional failure certificate. These are
useful dependencies, but their nominal types do not prove that the gate is sound,
that obligations are discharged, or that the returned statement equals the
requested statement. A dependent specification could index those obligations by
$q$; this diagram reports the existing signature.

\paragraph{Constructed example.}
The queue requests exponential decay for a finite connected undirected graph's symmetric Laplacian on
the mean-zero subspace for $\dot u=-Lu$, assuming a gap
$\lambda_1\geq c>0$. Deduction retains these premises in the Euclidean-norm
claim $\|u(t)\|_2\leq e^{-ct}\|u(0)\|_2$ for $t\geq0$.
Refutation checks the excluded constant mode and the gap assumption. A missing
spectral argument remains an open obligation; favorable testing cannot close it.
The gate must distinguish an accepted proof from a certificate explaining that
gap or quarantine.

Original ERFL routes accepted proof status to contribution cultivation, failure
certificates to conceptual exploration, and quarantine to stop. Native comments
preserve that policy as non-executable handoffs, including adaptation to a result
atlas. The native sequence supplies no local terminal acceptance transition.

% SPDX-FileCopyrightText: © 2026 Idiap Research Institute <contact@idiap.ch>
%
% SPDX-FileContributor: Andre Freitas <andre.freitas@idiap.ch>
% Neuro-symbolic AI Group
%
% SPDX-FileContributor: Danilo Gusicuma <danilo.gusicuma@idiap.ch>
% Neuro-symbolic AI Group
%
% SPDX-License-Identifier: MIT
% Generated from step-explanations.json
\paragraph{Step-by-step reconstruction.}
Prove Euclidean-norm decay for the symmetric combinatorial Laplacian of a finite connected undirected graph on mean-zero initial data, assuming its first positive eigenvalue is at least c greater than zero.
The following is a source-guided illustrative reconstruction. Internal procedures are prescriptions; the native declarations do not establish their execution.
\begin{enumerate}[leftmargin=*,label=\arabic*.]
\item \textbf{Construct the exact conditional theorem} (\texttt{d0}).
\begin{description}[style=standard,labelindent=0pt,leftmargin=5.5em,labelwidth=5em,labelsep=0.5em,align=left,font=\normalfont\bfseries,itemsep=0.15em,parsep=0pt,topsep=0.3em]
\item[Purpose] Derive the exact queued obligation.
\item[Input] item.premises: queued exact\_\allowbreak{}statement; connected-graph, subspace and gap premises
\item[Method] Order lemmas without circular dependence and propagate constants, norms and quantifiers through each inference.
\item[Check] Every inferential step must identify its premises and preserve the theorem's scope.
\item[Output] Exact claim; retained premises and undischarged formal obligations
\item[Next] Refutation receives this claim; a missing lemma remains visible rather than silently closed.
\item[Example] For du/dt=-Lu on the connected graph, expand the mean-zero vector in positive-eigenvalue modes. Each coefficient evolves by exp(-lambda t), giving Euclidean-norm decay at rate at least c.
\end{description}
\item \textbf{Search for a defeater} (\texttt{r0}).
\begin{description}[style=standard,labelindent=0pt,leftmargin=5.5em,labelwidth=5em,labelsep=0.5em,align=left,font=\normalfont\bfseries,itemsep=0.15em,parsep=0pt,topsep=0.3em]
\item[Purpose] Challenge edge cases without weakening the claim.
\item[Input] Exact claim; gap premise; open obligations
\item[Method] Test small, degenerate and adversarial cases; classify a failed step as a false statement, wrong quantifier, model mismatch or missing technique.
\item[Check] Favorable examples cannot substitute for a missing spectral argument.
\item[Output] Refutation assessment: scoped surviving theorem or localized defect
\item[Next] Calibration retains this finding with the exact claim; repair must name any changed premise.
\item[Example] The constant mode violates unrestricted decay but is excluded by the mean-zero premise. Without the gap premise, cycle families defeat a common positive rate.
\end{description}
\item \textbf{Calibrate the proof disposition} (\texttt{e0}).
\begin{description}[style=standard,labelindent=0pt,leftmargin=5.5em,labelwidth=5em,labelsep=0.5em,align=left,font=\normalfont\bfseries,itemsep=0.15em,parsep=0pt,topsep=0.3em]
\item[Purpose] Distinguish proof discharge from conditional or blocked work.
\item[Input] Exact theorem claim + its refutation assessment
\item[Method] Check the lemma chain and retained obligations; issue a bounded disposition and failure certificate when a bridge remains missing.
\item[Check] A proof gate cannot discharge an obligation merely because refutation found no counterexample.
\item[Output] Claim; finding; proof gate; optional failure certificate
\item[Next] Original policy routes accepted work or failure; the native return is no terminal proof acceptance act.
\item[Example] Accept only the explicitly justified conditional spectral estimate; if spectral decomposition has not been supplied, retain that proof debt in the disposition.
\end{description}
\end{enumerate}
\paragraph{Changed-input contrast.}
If the queue requests all initial data, a nonzero constant vector defeats positive-rate decay. The theorem must retain the mean-zero restriction; an assessment of that restricted theorem cannot certify the broader one.


% SPDX-FileCopyrightText: © 2026 Idiap Research Institute <contact@idiap.ch>
%
% SPDX-FileContributor: Andre Freitas <andre.freitas@idiap.ch>
% Neuro-symbolic AI Group
%
% SPDX-FileContributor: Danilo Gusicuma <danilo.gusicuma@idiap.ch>
% Neuro-symbolic AI Group
%
% SPDX-License-Identifier: MIT
\subsection{Target navigation: choose an obligation before doing research}
\label{app:granular-target-navigation-loop}
The linked dependency-cartographer skill supplies the conceptual procedure hidden
by this compact boundary. It reads the target's acceptance conditions against a
recursive dependency graph, distinguishes a missing premise from a bridge object,
route lemma, proof debt or stale route, and groups unresolved nodes into coherent
proof tasks. Priority is judged by how a node serves the target, not merely by
its local score. Similarity does not count as logical dependence, and a mentioned
lemma does not count as a discharged obligation. Cycle detection prevents a route
from supporting itself through a chain of references. The crucial inference is
from the \emph{structure of what is missing} to a justified next research question.
This is a prescribed interpretation procedure linked to the workflow, rather than
an algorithm encoded by the native projection signature.

Target navigation separates identifying what remains to be established from
establishing it. Write $T$ for \texttt{Unscoped\allowbreak Research\allowbreak Target}, $F$ for
\texttt{Target\allowbreak Dependency\allowbreak Frontier}, and $N$ for
\texttt{Target\allowbreak Navigation\allowbreak Disposition}. The native occurrences are an unresolved
projection $P_T$ and an evidential assessment $E_T$:
\[
\begin{aligned}
&P_T:T\to\operatorname{Generated}[F],\qquad
 E_T:F\to\operatorname{Generated}[N],\\
&f=P_T(t),\qquad n=E_T(f.\mathrm{value}),\\
&N=\{\mathrm{WorkReady},\mathrm{ReadyForReleasePromotion},\mathrm{Stalled}\}.
\end{aligned}
\]
The first occurrence, \texttt{dependency\_\allowbreak frontier}, changes an identified target
into a target with a referenced dependency frontier. It makes the unresolved
obligations available for selection, but adds no warrant for the target.
Crucially, the original \texttt{pre:\_\allowbreak } does not encode the linked procedure:
$P_T$ names a declared input/output boundary, not a recovered cartography
algorithm. The second occurrence, \texttt{target\_\allowbreak disposition}, assesses that
frontier to choose the next kind of work. The relation
\texttt{Dependency\allowbreak Frontier\allowbreak Enables\allowbreak Target\allowbreak Disposition} records the declared dependency
of the generated decision on the frontier; its generic source/target certificate
shape does not itself establish that the frontier is complete. The displayed
$\mathrm{value}$ selection is the native explicit wrapper adapter.

For a constructed example, suppose the target is a decay bound uniform over a
family of networks. A frontier can distinguish the proved conditional estimate
from the open obligation that the spectral gaps have a positive common lower
bound. Since that obligation remains open, \texttt{Work\allowbreak Ready} directs work to the
gap question; it cannot mean that uniform decay has been proved. This illustrates
why selecting an obligation is a distinct conceptual move from deriving a result.
The native function returns $n$. The original control policy instead switches
work-ready to \texttt{abducer\_\allowbreak work\_\allowbreak loop}, release-candidate to
\texttt{release\_\allowbreak promotion\_\allowbreak loop}, and stops on stall. Those transfers are not
executed by the native return statement, and the native declarations do not encode the skill-level
judgment that chooses among them.

% SPDX-FileCopyrightText: © 2026 Idiap Research Institute <contact@idiap.ch>
%
% SPDX-FileContributor: Andre Freitas <andre.freitas@idiap.ch>
% Neuro-symbolic AI Group
%
% SPDX-FileContributor: Danilo Gusicuma <danilo.gusicuma@idiap.ch>
% Neuro-symbolic AI Group
%
% SPDX-License-Identifier: MIT
% Generated from step-explanations.json
\paragraph{Step-by-step reconstruction.}
The target is a diffusion bound uniform across a graph family. A conditional decay theorem is available, but its common spectral-gap premise is open.
The following is a source-guided illustrative reconstruction. Internal procedures are prescriptions; the native declarations do not establish their execution.
\begin{enumerate}[leftmargin=*,label=\arabic*.]
\item \textbf{Construct the dependency frontier} (\texttt{dependency\_\allowbreak{}frontier}).
\begin{description}[style=standard,labelindent=0pt,leftmargin=5.5em,labelwidth=5em,labelsep=0.5em,align=left,font=\normalfont\bfseries,itemsep=0.15em,parsep=0pt,topsep=0.3em]
\item[Purpose] Locate the target’s remaining proof obligations.
\item[Input] Target acceptance conditions; recursive dependency graph
\item[Method] The native producer is unresolved. The linked cartographer prescribes classifying missing premises, bridge objects, route lemmas and proof debt, while checking cycles.
\item[Check] A dependency must support the target logically; similarity or a mentioned lemma does not discharge it.
\item[Output] Intended frontier: open uniform-gap proof obligation
\item[Next] Disposition assessment requires this frontier; the signature does not establish that cartography was performed.
\item[Example] Separate the available conditional decay theorem from the unproved assertion that every graph has gap at least the same c greater than zero.
\end{description}
\item \textbf{Select a justified obligation} (\texttt{target\_\allowbreak{}disposition}).
\begin{description}[style=standard,labelindent=0pt,leftmargin=5.5em,labelwidth=5em,labelsep=0.5em,align=left,font=\normalfont\bfseries,itemsep=0.15em,parsep=0pt,topsep=0.3em]
\item[Purpose] Select work that controls target completion.
\item[Input] Dependency frontier, including unresolved support
\item[Method] Assess which open node controls the target and distinguish work readiness from release readiness or stall.
\item[Check] An open necessary premise prevents treating the target as complete.
\item[Output] Work disposition: common-gap task and justification
\item[Next] Original policy routes onward; the native return neither proves the premise nor executes that route.
\item[Example] Because uniform decay depends on the open common-gap premise, select that premise for investigation rather than re-deriving the conditional theorem.
\end{description}
\end{enumerate}
\paragraph{Changed-input contrast.}
If the common-gap premise is discharged for the target family, navigation should reassess the remaining acceptance obligations. It must not continue selecting the same already closed gap merely because its node is prominent.


\subsection{Release promotion: conserve the scope of the evidence}
\label{app:granular-release-promotion-loop}
The release-promoter skill prescribes a comparative interpretation of evidence:
separate robust support from missing or fragile support, inspect whether the
active route still addresses the target, and determine whether the chosen
mathematical machinery has the objects it needs. This last question can change the
research decision even when every report validates: endpoint spectral machinery
cannot replace an absent construction or transfer object. A route that repeatedly
hits the same blocker must explain what was learned and what would trigger a
pivot. Promotion consequently depends on whether the proposed claim survives
these comparisons; repair can change the target scope or conceptual route rather
than merely regenerate a report. These are skill-level prescriptions associated
with the release contract, not extra native occurrences or evidence of execution.

A release decision bounds the result supported by a research iteration.
Here $V$ denotes \texttt{Promotion\allowbreak Evidence\allowbreak Packet}, $Q$ denote
\texttt{Claim\allowbreak Bounded\allowbreak Release\allowbreak Proposal}, and $D$ denote
\texttt{Release\allowbreak Promotion\allowbreak Disposition}. Its typed boundary is
\[
\begin{aligned}
&P_R:V\to\operatorname{Generated}[Q],\qquad
 E_R:Q\to\operatorname{Generated}[D],\\
&q=P_R(v),\qquad d=E_R(q.\mathrm{value}),\\
&D=\{\mathrm{Promoted},\mathrm{RepairTarget},
       \mathrm{RepairConcept},\mathrm{Rejected}\}.
\end{aligned}
\]
At \texttt{release\_\allowbreak proposal}, the evidence packet's run and evidence references
are projected into a proposal with run-identity and proposal-reference fields.
Preserving run identity is intended, but equality of the two string fields is
not enforced by these nominal types. Conceptually this is the move from accumulated evidence to a bounded
claim about it. The source's \texttt{pre:\_\allowbreak } leaves proposal construction open;
\texttt{Unresolved\allowbreak Release\allowbreak Proposal\allowbreak Projection} must not be read as an implemented
claim-selection algorithm. At \texttt{promotion\_\allowbreak disposition}, evidential
assessment asks whether that proposal is supportable, needs a target or conceptual
repair, or should be rejected. \texttt{Release\allowbreak Proposal\allowbreak Enables\allowbreak Promotion\allowbreak Disposition}
licenses the declared proposal-to-assessment dependency. It supplies neither an
independent proof nor a rule that reference-bearing proposals must pass. There is
no native multi-input join: the assessment receives the unwrapped proposal, whose
reference must carry the evidence boundary.

For a constructed example, a Yang--Mills route has a finite-regulator correlation
estimate and proposes promoting it as evidence that its continuum construction is
complete. The release comparison exposes an absent same-theory transfer and
nonuniform error control. Repairing the target can license a finite-result claim;
repairing the concept instead asks for the missing transfer object. Passing the
workflow validators cannot decide between those scientific scopes. The skill's
machinery assessment makes the difference explicit: a spectral conclusion cannot
be obtained by invoking endpoint machinery before its analytic preconditions are
available.
The original policy accepts a promoted outcome, routes target repair to target
navigation and conceptual repair to conceptual hypothesis exploration, and stops
on rejection. The native function returns the wrapped disposition and contains
none of those switches. Thus the workflow's contribution is scope calibration,
not an additional experiment or an automatic truth certificate.

% SPDX-FileCopyrightText: © 2026 Idiap Research Institute <contact@idiap.ch>
%
% SPDX-FileContributor: Andre Freitas <andre.freitas@idiap.ch>
% Neuro-symbolic AI Group
%
% SPDX-FileContributor: Danilo Gusicuma <danilo.gusicuma@idiap.ch>
% Neuro-symbolic AI Group
%
% SPDX-License-Identifier: MIT
% Generated from step-explanations.json
\paragraph{Step-by-step reconstruction.}
A release packet contains a finite-regulator correlation estimate but proposes a continuum conclusion without a same-theory transfer or uniform error bound.
The following is a source-guided illustrative reconstruction. Internal procedures are prescriptions; the native declarations do not establish their execution.
\begin{enumerate}[leftmargin=*,label=\arabic*.]
\item \textbf{Construct a bounded proposal} (\texttt{release\_\allowbreak{}proposal}).
\begin{description}[style=standard,labelindent=0pt,leftmargin=5.5em,labelwidth=5em,labelsep=0.5em,align=left,font=\normalfont\bfseries,itemsep=0.15em,parsep=0pt,topsep=0.3em]
\item[Purpose] Bound a release claim by accumulated support.
\item[Input] Iteration evidence packet; target identity; support references
\item[Method] The native construction is unresolved. The linked promoter prescribes separating robust support from fragile support and exposing target drift and missing machinery prerequisites.
\item[Check] The proposed scope cannot exceed the regime actually supported by evidence.
\item[Output] Intended proposal: supported finite estimate; unsupported continuum extension
\item[Next] Promotion assessment consumes the proposal; nominal run references do not enforce identity preservation.
\item[Example] Write the finite-regulator estimate as supported within its assumptions and list continuum passage and uniform error control as open obligations.
\end{description}
\item \textbf{Assess promotion or repair} (\texttt{promotion\_\allowbreak{}disposition}).
\begin{description}[style=standard,labelindent=0pt,leftmargin=5.5em,labelwidth=5em,labelsep=0.5em,align=left,font=\normalfont\bfseries,itemsep=0.15em,parsep=0pt,topsep=0.3em]
\item[Purpose] Assess promotion against the evidence boundary.
\item[Input] Release proposal; underlying support references
\item[Method] Compare claim scope with missing construction and transfer objects; distinguish target repair, conceptual repair and rejection.
\item[Check] Endpoint spectral machinery cannot replace an absent analytic transfer premise.
\item[Output] Promotion disposition: bounded promotion or specified repair
\item[Next] Original policy may promote, route repair or stop; native output alone performs none of these effects.
\item[Example] Restrict the release to the finite estimate or request the missing transfer object; do not accept the continuum announcement from the same packet.
\end{description}
\end{enumerate}
\paragraph{Changed-input contrast.}
Restricting the proposal to the evidenced finite-regulator result can make a bounded claim supportable. Adding more fixed-regulator examples cannot by itself repair the missing uniform transfer.


\subsection{Timeline-to-paper: translate a reasoning record without upgrading it}
\label{app:granular-timeline-paper-independent-workflow}
The linked argument-architect and failed-route-interpreter skills explain how
the record becomes a scientific argument. The architect chooses the paper's form
from the evidence: theorem, worked calculation, obstruction, method, negative
result or positioning note. It identifies the smallest defensible contribution
and maps it to the evidence, novelty boundary and remaining validation burden.
The failure interpreter compresses a route into its original idea, the pressure
that defeated it, the learned invariant and the narrowed object that survives.
Their conceptual product is an explanation of why the current claim is the right
claim, rather than a chronological list of attempts. Section responsibilities
then distribute that argument across the paper. The skill procedures specify this
reasoning even though the native claim-set projection does not implement it;
manifest membership alone supplies no execution order between the linked skills.

Paper construction has an epistemic boundary: reorganizing a reasoning record
must not turn its open conjectures into established results. Let $L$ be
\texttt{Timeline\allowbreak Reasoning\allowbreak Outputs}, $C$ be
\texttt{Bounded\allowbreak Manuscript\allowbreak Claim\allowbreak Set}, and $D$ be
\texttt{Timeline\allowbreak Paper\allowbreak Disposition}. The two declared occurrences give
\[
\begin{aligned}
&P_P:L\to\operatorname{Generated}[C],\qquad
 E_P:C\to\operatorname{Generated}[D],\\
&c=P_P(\ell),\qquad d=E_P(c.\mathrm{value}),\\
&D=\{\mathrm{PaperReady},\mathrm{ClaimBoundaryFailure},\mathrm{Stalled}\}.
\end{aligned}
\]
The \texttt{bounded\_\allowbreak claim\_\allowbreak set} occurrence translates a timeline reference into
a source reference and claim-ledger reference. This establishes presentation
input, not new evidence. Its unresolved projection corresponds to the original
\texttt{pre:\_\allowbreak }; the declarations do not specify how competing claims are
selected, reconciled, or written. The \texttt{paper\_\allowbreak disposition} occurrence
assesses whether this bounded set supports paper assembly.
\texttt{Claim\allowbreak Set\allowbreak Enables\allowbreak Paper\allowbreak Disposition} records the assessment's dependence on
that claim set, with explicit unwrapping at $c.\mathrm{value}$. It does not assert
that the ledger faithfully reflects its source merely because both references
exist. Neither types nor topology specify an internal drafting sequence.

Consider a constructed diffusion study whose timeline contains a conditional
decay derivation, an unsuccessful attempt to generalize it, and the cycle-graph
counterexample to a size-independent spectral gap. A faithful claim set separates
these three statuses. The paper can explain why the generalization fails while
retaining the conditional theorem. If a draft instead states uniform decay for
all connected graphs, \texttt{Claim\allowbreak Boundary\allowbreak Failure} is an appropriate assessment:
the presentation has dropped precisely the assumption challenged by the record.
The outcome is a corrected scientific claim boundary, not merely a formatted
manuscript. The original workflow accepts paper-ready and stops on claim-boundary
failure or stall. The native function only returns $d$; its readiness value is
not itself a publication act. This distinction matters because readable exposition
can improve understanding without increasing the evidential status of any claim.

% SPDX-FileCopyrightText: © 2026 Idiap Research Institute <contact@idiap.ch>
%
% SPDX-FileContributor: Andre Freitas <andre.freitas@idiap.ch>
% Neuro-symbolic AI Group
%
% SPDX-FileContributor: Danilo Gusicuma <danilo.gusicuma@idiap.ch>
% Neuro-symbolic AI Group
%
% SPDX-License-Identifier: MIT
% Generated from step-explanations.json
\paragraph{Step-by-step reconstruction.}
A diffusion timeline contains a conditional decay derivation, a failed uniform generalization and a cycle-graph counterexample to a common spectral gap.
The following is a source-guided illustrative reconstruction. Internal procedures are prescriptions; the native declarations do not establish their execution.
\begin{enumerate}[leftmargin=*,label=\arabic*.]
\item \textbf{Form the defensible claim set} (\texttt{bounded\_\allowbreak{}claim\_\allowbreak{}set}).
\begin{description}[style=standard,labelindent=0pt,leftmargin=5.5em,labelwidth=5em,labelsep=0.5em,align=left,font=\normalfont\bfseries,itemsep=0.15em,parsep=0pt,topsep=0.3em]
\item[Purpose] Recover a supported argument from the research record.
\item[Input] Timeline claims; premises; counterexamples; unresolved attempts
\item[Method] The native projection is unresolved. The argument architect selects the smallest supported contribution; failure interpretation extracts the defeated idea, learned obstruction and surviving statement.
\item[Check] The paper mode must follow the evidence rather than the ambition of the original target.
\item[Output] Intended claim ledger: conditional theorem; refuted generalization
\item[Next] Paper assessment consumes this bounded ledger; the native boundary supplies no automatic drafting algorithm.
\item[Example] Preserve the conditional spectral estimate, report cycles as the obstruction to uniformity, and explain why the unrestricted claim was abandoned.
\end{description}
\item \textbf{Check presentation readiness} (\texttt{paper\_\allowbreak{}disposition}).
\begin{description}[style=standard,labelindent=0pt,leftmargin=5.5em,labelwidth=5em,labelsep=0.5em,align=left,font=\normalfont\bfseries,itemsep=0.15em,parsep=0pt,topsep=0.3em]
\item[Purpose] Check that manuscript organization preserves claim scope.
\item[Input] Bounded claim set; source references
\item[Method] Compare proposed manuscript claims with the premises and failure record, locating any omitted assumption or contradictory status.
\item[Check] A claim-ledger reference is insufficient unless its content faithfully represents the record.
\item[Output] Paper disposition: readiness or specified claim-boundary failure
\item[Next] Original policy accepts readiness or stops; native return is not publication.
\item[Example] A draft asserting uniform decay for all connected graphs fails because it removes the very common-gap premise challenged by the cycle example.
\end{description}
\end{enumerate}
\paragraph{Changed-input contrast.}
Removing the counterexample from the narrative while retaining the broad claim breaks fidelity to the record. Restoring the failed premise and scoped conditional theorem repairs the claim boundary without creating new evidence.


\subsection{Yang--Mills overlay: keep representation and continuum claims apart}
\label{app:granular-ym-representational-dynamics-overlay}
The linked Yang--Mills skills describe a much richer scientific program behind
this audit boundary. Observable-core construction asks whether the chosen
observables separate states and generate enough physical vectors to see every
relevant spectral sector. Representation construction independently specifies
source and target carriers, observables, states and dynamics, then asks what an
observable pullback and state pushforward preserve. Exact equivalence is
separated from coarse-graining and approximation with explicit error dependence.
Continuum coherence asks whether regulated correlations converge in a topology
strong enough to preserve positivity and spectral estimates, with justified limit
orders and theory identity. Rigidity transfer then asks whether a uniform decay
estimate extends from a dense physical class to spectral exclusion. These are
prescribed mathematical burdens, not solved theorems or an execution sequence
inferred from the skill list. The native two-stage audit compresses their results;
it does not make the underlying conceptual questions unknown.

This overlay tests whether a representation candidate stays within the declared
continuum boundary. Write $Y$ for \texttt{Yang\allowbreak Mills\allowbreak Target\allowbreak Context}, $A$ for
\texttt{Continuum\allowbreak Boundary\allowbreak Audit\allowbreak Input}, and $D$ for
\texttt{Yang\allowbreak Mills\allowbreak Overlay\allowbreak Disposition}:
\[
\begin{aligned}
&P_Y:Y\to\operatorname{Generated}[A],\qquad
 E_Y:A\to\operatorname{Generated}[D],\\
&a=P_Y(y),\qquad d=E_Y(a.\mathrm{value}),\\
&D=\{\mathrm{Accepted},\mathrm{Quarantined},\mathrm{Stalled}\}.
\end{aligned}
\]
At \texttt{continuum\_\allowbreak audit\_\allowbreak input}, the projection supplies target-identity and
continuum-boundary fields together with a representation-candidate reference.
Those fields can retain the input context, but preserving their values is an
obligation not enforced by the nominal string types.
Its conceptual purpose is to put a candidate beside the boundary it must respect.
It supplies no bridge from a finite representation to a continuum assertion.
The original \texttt{pre:\_\allowbreak } leaves that preparation internally unspecified;
one cannot infer literature search, simulation, or theorem construction as extra
native occurrences. At \texttt{overlay\_\allowbreak disposition}, evidential assessment
judges the prepared candidate against the boundary.
\texttt{Audit\allowbreak Input\allowbreak Enables\allowbreak Overlay\allowbreak Disposition} connects this audit input to its
generated assessment, without certifying that the boundary has been crossed or
that an accepted representation proves the scientific target.

For a constructed illustration, suppose a selected observable has exponentially
decaying correlations. Add an orthogonal sector unseen by that observable whose
positive energies approach zero. The selected correlation is unchanged, yet a
global spectral-gap claim is defeated. Observable-core reasoning identifies the
missing density obligation; rigidity transfer restricts the conclusion to the
visible sector until that obligation is met. This countermodel distinguishes
``the measurement decays'' from ``all physical excitations are gapped.'' It is a
test of the prescribed inference, not a proposed Yang--Mills construction.
The original ERFL additionally queries the connected knowledge base during
preparation and persists the audit. These effects have no corresponding native
query or persistence occurrence in this two-stage projection. Original accepted,
quarantined, and stalled outcomes lead respectively to accept or stop, whereas the
native function returns $d$. Thus acceptance concerns an overlay disposition;
it must not silently acquire the meaning of a continuum existence proof.

% SPDX-FileCopyrightText: © 2026 Idiap Research Institute <contact@idiap.ch>
%
% SPDX-FileContributor: Andre Freitas <andre.freitas@idiap.ch>
% Neuro-symbolic AI Group
%
% SPDX-FileContributor: Danilo Gusicuma <danilo.gusicuma@idiap.ch>
% Neuro-symbolic AI Group
%
% SPDX-License-Identifier: MIT
% Generated from step-explanations.json
\paragraph{Step-by-step reconstruction.}
A selected observable shows exponential decay, while the proposed global spectral conclusion has no proof that the observable-generated sector is dense.
The following is a source-guided illustrative reconstruction. Internal procedures are prescriptions; the native declarations do not establish their execution.
\begin{enumerate}[leftmargin=*,label=\arabic*.]
\item \textbf{Prepare the continuum-boundary audit} (\texttt{continuum\_\allowbreak{}audit\_\allowbreak{}input}).
\begin{description}[style=standard,labelindent=0pt,leftmargin=5.5em,labelwidth=5em,labelsep=0.5em,align=left,font=\normalfont\bfseries,itemsep=0.15em,parsep=0pt,topsep=0.3em]
\item[Purpose] Bind the representation to its physical claim boundary.
\item[Input] Target context; representation proposal; continuum boundary
\item[Method] The native preparation is unresolved. Linked procedures specify carriers, observable maps, state transport, positivity, limit order and observable completeness as separate burdens.
\item[Check] Finite or sectorial evidence must remain distinguished from same-theory continuum and full-spectrum claims.
\item[Output] Intended audit input: target identity; boundary; candidate reference
\item[Next] Overlay assessment consumes it; query and persistence effects in original policy are not native occurrences.
\item[Example] Record decay for the visible observable family and mark density of its generated physical sector as an unproved assumption.
\end{description}
\item \textbf{Test the global inference} (\texttt{overlay\_\allowbreak{}disposition}).
\begin{description}[style=standard,labelindent=0pt,leftmargin=5.5em,labelwidth=5em,labelsep=0.5em,align=left,font=\normalfont\bfseries,itemsep=0.15em,parsep=0pt,topsep=0.3em]
\item[Purpose] Identify unsupported physical extrapolation.
\item[Input] Prepared candidate; retained continuum limits
\item[Method] Seek a sector invisible to the chosen observables and inspect whether the proposed transfer preserves the required structure.
\item[Check] Unchanged visible correlations cannot exclude low energies in an unobserved sector.
\item[Output] Overlay disposition: limited claim or quarantined global inference
\item[Next] Native output returns an overlay disposition; accepted status is not a continuum existence or mass-gap theorem.
\item[Example] In an abstract spectral calibration, adjoin an orthogonal sector with positive energies approaching zero. Visible decay is unchanged; this exposes the missing density inference without constructing a Yang--Mills counterexample.
\end{description}
\end{enumerate}
\paragraph{Changed-input contrast.}
If density and the required uniform extension estimates were independently established, the hidden-sector objection would need reassessment. Their mere appearance as named fields in an audit packet does not establish them.


\subsection{Deep review: challenge and repair the same epistemic state}
\label{app:granular-deep-review-and-reorientation}
The linked review skill makes disagreement itself an object of reasoning.
Reviewers inspect implemented code and persisted state through mathematical,
epistemic, conceptual, discovery and design lenses; they separate a mathematical
defect from an implementation defect and compare where the diagnoses converge or
conflict. Reflection then asks what the search changed in the understanding of
the problem: a new perspective, a sharper obstruction, persistent hardness or a
regression. A more elaborate artifact is not automatically a deeper result.
The named contract prescribes state formalization, advisory reviews, meta-review
and repair governance; this procedural ordering is a different layer from the
native parallel challenge/repair topology. Together they explain how a failure
becomes a revised research question while preserving unresolved disagreement,
rather than averaging reviewer opinions into apparent confirmation.

Deep review separates detecting why a research direction fails from proposing how
it might recover. Let $B,S,F,P,J,D$ abbreviate respectively
\texttt{Review\allowbreak Evidence\allowbreak Bundle}, \texttt{Formalized\allowbreak Review\allowbreak State},
\texttt{Review\allowbreak Defeater\allowbreak Assessment}, \texttt{Reorientation\allowbreak Repair\allowbreak Options},
\texttt{Reorientation\allowbreak Decision\allowbreak Input}, and \texttt{Reorientation\allowbreak Decision\allowbreak Assessment}.
\[
\begin{aligned}
&G:B\to\operatorname{Generated}[S],\quad
 R:S\to\operatorname{Generated}[F],\quad
 A:S\to\operatorname{Generated}[P],\\
&E:J\to\operatorname{Generated}[D],\qquad s=G(b),\\
&\left.\begin{aligned} f&=R(s.\mathrm{value})\\
                        p&=A(s.\mathrm{value})\end{aligned}\right\}
 \text{ declared parallel},\\
&j=\{\mathrm{defeater\_assessment}=f.\mathrm{value},\;
       \mathrm{repair\_options}=p.\mathrm{value}\},\qquad d=E(j).
\end{aligned}
\]
\texttt{review\_\allowbreak state} grounds the accumulated evidence in one formalized state.
\texttt{defeater\_\allowbreak assessment} attacks that state to identify bounded defects,
while \texttt{repair\_\allowbreak options} abductively proposes changes from the same state.
The two branches are complementary, not competing votes.
\texttt{Grounds\allowbreak Reorientation\allowbreak Alternatives} annotates both dependencies and makes
explicit that their point of departure is shared.
\texttt{reorientation\_\allowbreak decision} then calibrates the response using both named
fields of $J$. \texttt{Calibrates\allowbreak Reorientation\allowbreak Decision} records this joined
assessment boundary. The relation records express declared dependencies, not a
proof that a suggested repair resolves the defeater. Importantly, the repair
branch does not consume $f$: a repair proposed in parallel still has to be judged
against the defeater at the join.

For a constructed example, one review sees a promising correlation-decay
calculation, while another finds that the chosen observables miss a possible
low-energy sector. The disagreement is resolved by identifying its object: the
calculation may be valid for a visible sector while the global conclusion lacks
density. Reflection records a sharpening of the obstruction, not a completed gap
proof. Repair must construct a sufficiently rich observable class or restrict the
claim; simply collecting more decay curves would not address the defeater.
The output pairs a decision explanation with repair-assigned,
conceptual-reorientation-required, or stalled disposition. Original control routes
the first two to target navigation or conceptual hypothesis exploration and stops
on stall. Native reasoning returns the wrapped assessment; declared parallelism
and routing do not establish an observed execution.

% SPDX-FileCopyrightText: © 2026 Idiap Research Institute <contact@idiap.ch>
%
% SPDX-FileContributor: Andre Freitas <andre.freitas@idiap.ch>
% Neuro-symbolic AI Group
%
% SPDX-FileContributor: Danilo Gusicuma <danilo.gusicuma@idiap.ch>
% Neuro-symbolic AI Group
%
% SPDX-License-Identifier: MIT
% Generated from step-explanations.json
\paragraph{Step-by-step reconstruction.}
Reviewers examine a correlation-decay argument whose calculation concerns a visible sector but whose conclusion concerns the full physical spectrum.
The following is a source-guided illustrative reconstruction. Internal procedures are prescriptions; the native declarations do not establish their execution.
\begin{enumerate}[leftmargin=*,label=\arabic*.]
\item \textbf{Ground the reviewed state} (\texttt{review\_\allowbreak{}state}).
\begin{description}[style=standard,labelindent=0pt,leftmargin=5.5em,labelwidth=5em,labelsep=0.5em,align=left,font=\normalfont\bfseries,itemsep=0.15em,parsep=0pt,topsep=0.3em]
\item[Purpose] Ground criticism and repair in the same state.
\item[Input] Evidence bundle; implemented argument; persisted claims
\item[Method] Separate mathematical results, implementation behavior, assumptions and open obligations before interpreting progress.
\item[Check] Disagreement about scope must remain explicit rather than averaged into confidence.
\item[Output] Formalized review state: visible decay; unsupported global extension
\item[Next] Both parallel branches consume this same state; their coexistence gives no independence guarantee.
\item[Example] Write the state as a valid-looking visible-sector calculation plus the additional, unsupported claim that every physical excitation is gapped.
\end{description}
\item \textbf{Localize the defeater} (\texttt{defeater\_\allowbreak{}assessment}).
\begin{description}[style=standard,labelindent=0pt,leftmargin=5.5em,labelwidth=5em,labelsep=0.5em,align=left,font=\normalfont\bfseries,itemsep=0.15em,parsep=0pt,topsep=0.3em]
\item[Purpose] Locate the unsupported implication.
\item[Input] Grounded review state, shared with repair formation
\item[Method] Challenge quantifiers and representation coverage, separating a defect in the inference from a defect in implementation.
\item[Check] A useful diagnosis names the missing burden and the claim it affects.
\item[Output] Defeater assessment: missing observable-density premise
\item[Next] The join later compares this finding with repairs; the repair branch does not consume it directly.
\item[Example] An abstract direct-sum spectral countermodel preserves the reported correlations while adding energies approaching zero. It defeats the unrestricted correlation-to-gap inference and identifies the missing completeness premise; it is not a constructed Yang--Mills theory.
\end{description}
\item \textbf{Propose recovery options} (\texttt{repair\_\allowbreak{}options}).
\begin{description}[style=standard,labelindent=0pt,leftmargin=5.5em,labelwidth=5em,labelsep=0.5em,align=left,font=\normalfont\bfseries,itemsep=0.15em,parsep=0pt,topsep=0.3em]
\item[Purpose] Propose repairs that preserve useful content.
\item[Input] Formalized review state; no parallel refutation output
\item[Method] Propose a changed representation, missing bridge or narrower claim, explaining what each option would preserve.
\item[Check] A repair proposal must alter the defective dependency rather than relabel the current result.
\item[Output] Repair alternatives: richer observables or sectorial theorem
\item[Next] Final calibration must test these options against the separately produced defeater.
\item[Example] From the visible/global mismatch, propose either constructing a dense observable-generated subspace or explicitly restricting the conclusion to the visible sector.
\end{description}
\item \textbf{Judge repairs against criticism} (\texttt{reorientation\_\allowbreak{}decision}).
\begin{description}[style=standard,labelindent=0pt,leftmargin=5.5em,labelwidth=5em,labelsep=0.5em,align=left,font=\normalfont\bfseries,itemsep=0.15em,parsep=0pt,topsep=0.3em]
\item[Purpose] Match repairs to diagnosed defects.
\item[Input] Defeater assessment + parallel repair options
\item[Method] Check which proposal addresses the missing implication and record whether understanding sharpened, stalled or regressed.
\item[Check] A proposed repair does not resolve a defect merely by appearing in the same packet.
\item[Output] Reorientation assessment: repair assignment or conceptual reorientation
\item[Next] Original policy routes assignments; native return does not execute repair or prove density.
\item[Example] The narrower claim resolves the overstatement; the richer observable construction remains a research task. More curves without new sectors do not answer the defeater.
\end{description}
\end{enumerate}
\paragraph{Changed-input contrast.}
If a proposed repair only adds more decay curves from the same incomplete observable family, the join should reject it as irrelevant to density. A narrower sectorial claim changes the conclusion honestly.


\subsection{Contribution cultivation: test what a result enables}
\label{app:granular-scientific-contribution-cultivation}
The linked contribution skills separate finding a useful object from judging
its value and its effect on the target. Detection asks whether a subresult can
stand alone as a bridge, obstruction, reusable method or explanatory reframing,
then tests the smallest proposed claim against prior art, hidden premises and
proxy drift. Value assessment compares what would remain unknown without the
result, distinguishing exploration payoff from paper readiness. Alignment asks
whether the source object can feed the original target through a stated transfer;
enablement compares the precondition before the result with the postcondition
actually consumed by a next proof, falsifier or route decision. A beautiful local
calculation can therefore be valuable yet non-enabling, while a negative result
can enable a decisive pivot. The named contract explicitly separates these
assessments; they refine the meaning of the three native occurrences rather than
adding undocumented operations to that topology.

A supported result does not determine its scientific contribution by itself.
Contribution cultivation proposes an interpretation, challenges it, and assesses
readiness without inflating novelty. Let $S,H,T,J,Q$ abbreviate
\texttt{Result\allowbreak Atlas\allowbreak Slice}, \texttt{Scientific\allowbreak Contribution\allowbreak Hypothesis},
\texttt{Contribution\allowbreak Support\allowbreak Assessment}, \texttt{Contribution\allowbreak Readiness\allowbreak Input},
and \texttt{Contribution\allowbreak Readiness\allowbreak Assessment}:
\[
\begin{aligned}
&A:S\to\operatorname{Generated}[H],\qquad
 R:H\to\operatorname{Generated}[T],\qquad
 E:J\to\operatorname{Generated}[Q],\\
&h=A(s),\qquad t=R(h.\mathrm{value}),\\
&j=\{\mathrm{contribution\_hypothesis}=h.\mathrm{value},\;
       \mathrm{support\_assessment}=t.\mathrm{value}\},\qquad q=E(j).
\end{aligned}
\]
Occurrence \texttt{a0} abductively interprets the result as a candidate
contribution. Occurrence \texttt{r0} tests whether the result supports that
interpretation, under \texttt{Tests\allowbreak Contribution\allowbreak Hypothesis}. Occurrence
\texttt{e0} retains both interpretation and assessment at the named join and
calibrates readiness under \texttt{Calibrates\allowbreak Contribution\allowbreak Readiness}. The first
relation marks a test dependency; the second marks a calibration dependency.
Neither supplies missing evidence merely by being declared. The output contains
a disposition together with \texttt{Contribution\allowbreak Value\allowbreak Assessment},
\texttt{Main\allowbreak Problem\allowbreak Alignment\allowbreak Audit}, and \texttt{Enablement\allowbreak Audit}. These are returned
components, not three additional declared operation occurrences. Notice that
$q.\mathrm{value}$ unwraps \texttt{Generated}, whereas the assessment's own
\texttt{value} field denotes the contribution-value audit.

For a constructed example, cycle graphs show that connectedness alone cannot
support a spectral-gap lower bound uniform in graph size. An initial contribution
hypothesis might claim a general obstruction to fast diffusion. Testing rejects
that breadth: the result obstructs the uniform claim for the stated family and
assumptions. A defensible contribution can instead identify the missing structural
condition, align it with the uniform-decay problem, and enable a sharper search
for sufficient assumptions. Readiness still requires support for significance;
this illustration establishes neither novelty nor publication merit.
Original control stops when \texttt{r0} is refuted, routes paper-ready to
\texttt{manuscript\_\allowbreak result\_\allowbreak elaboration}, and stops on stall. The native sequence
always displays \texttt{e0} after \texttt{r0} and returns its wrapped assessment;
its handoff and stop comments are explicitly non-executable. A refuted
interpretation therefore must not be narrated as necessarily reaching the native
full-path calibration in the original policy.

% SPDX-FileCopyrightText: © 2026 Idiap Research Institute <contact@idiap.ch>
%
% SPDX-FileContributor: Andre Freitas <andre.freitas@idiap.ch>
% Neuro-symbolic AI Group
%
% SPDX-FileContributor: Danilo Gusicuma <danilo.gusicuma@idiap.ch>
% Neuro-symbolic AI Group
%
% SPDX-License-Identifier: MIT
% Generated from step-explanations.json
\paragraph{Step-by-step reconstruction.}
A completed cycle-graph example defeats a size-independent spectral gap from connectedness alone. The task is to identify what this negative result enables.
The following is a source-guided illustrative reconstruction. Internal procedures are prescriptions; the native declarations do not establish their execution.
\begin{enumerate}[leftmargin=*,label=\arabic*.]
\item \textbf{Interpret a possible contribution} (\texttt{a0}).
\begin{description}[style=standard,labelindent=0pt,leftmargin=5.5em,labelwidth=5em,labelsep=0.5em,align=left,font=\normalfont\bfseries,itemsep=0.15em,parsep=0pt,topsep=0.3em]
\item[Purpose] Identify the bounded result’s scientific use.
\item[Input] Result-atlas slice; assumptions; failure evidence
\item[Method] Propose the smallest standalone bridge, obstruction or reframing and identify who could use it.
\item[Check] A contribution hypothesis must not inherit the breadth of the original unsolved target.
\item[Output] Contribution hypothesis: obstruction to the insufficient premise
\item[Next] Refutation tests this interpretation, not merely the underlying calculation.
\item[Example] Propose that the cycle example shows why connectedness is too weak for the intended uniform diffusion bound and motivates a stronger structural condition.
\end{description}
\item \textbf{Challenge support and significance} (\texttt{r0}).
\begin{description}[style=standard,labelindent=0pt,leftmargin=5.5em,labelwidth=5em,labelsep=0.5em,align=left,font=\normalfont\bfseries,itemsep=0.15em,parsep=0pt,topsep=0.3em]
\item[Purpose] Test contribution scope and novelty support.
\item[Input] Proposed obstruction interpretation; stated scope
\item[Method] Test hidden premises, proxy drift and prior-art boundaries; compare what the result actually rules out.
\item[Check] Neither one slow family nor a true calculation establishes a general impossibility of fast diffusion.
\item[Output] Support assessment: bounded obstruction; unestablished breadth and novelty
\item[Next] Final readiness joins this finding with the same contribution hypothesis.
\item[Example] Reject the claim that fast diffusion is impossible in general; retain only failure of connectedness to imply a uniform gap over the specified family.
\end{description}
\item \textbf{Assess value and enablement} (\texttt{e0}).
\begin{description}[style=standard,labelindent=0pt,leftmargin=5.5em,labelwidth=5em,labelsep=0.5em,align=left,font=\normalfont\bfseries,itemsep=0.15em,parsep=0pt,topsep=0.3em]
\item[Purpose] Assess target value and communication readiness.
\item[Input] Contribution hypothesis + its support assessment
\item[Method] Separate conceptual value, target alignment and actual downstream enablement; compare the prior missing condition with the next task's requirements.
\item[Check] Exploration payoff cannot by itself establish paper readiness or novelty.
\item[Output] Readiness record: useful obstruction; open merit burden
\item[Next] Original policy can route to elaboration or stop; returned audit components are not extra native occurrences.
\item[Example] The result enables replacing an inadequate premise by a search for sufficient structural assumptions; any claim of publishable novelty still requires separate evidence.
\end{description}
\end{enumerate}
\paragraph{Changed-input contrast.}
If the result merely restates a known example without a new warranted use, its mathematical validity does not establish novelty or publication merit. Enablement needs a specific downstream change.



% SPDX-FileCopyrightText: © 2026 Idiap Research Institute <contact@idiap.ch>
%
% SPDX-FileContributor: Andre Freitas <andre.freitas@idiap.ch>
% Neuro-symbolic AI Group
%
% SPDX-FileContributor: Danilo Gusicuma <danilo.gusicuma@idiap.ch>
% Neuro-symbolic AI Group
%
% SPDX-License-Identifier: MIT
% Source-bound conceptual exposition; aliases are local to each subsection.
\subsection{Cohere Fast: deciding whether a scientific target is usable}
\label{app:granular-cohere_fast}
\texttt{cohere\_fast} separates identifying a scientific target from deciding whether inquiry can responsibly begin. Its two occurrences do not yet explain or prove the target. The first, \texttt{pre:\_} in the original description, projects a domain target into a scope activation. The native name, \texttt{Unresolved\allowbreak Cohere\allowbreak Scope\allowbreak Projection}, preserves this missing mechanism rather than classifying it as abduction. Intake calibration then assesses that activation and returns a readiness disposition. The relation \texttt{Scope\allowbreak Activation\allowbreak Enables\allowbreak Cohere\allowbreak Intake} expresses a prerequisite: an identifiable activation enables assessment, but provides no evidence that the target claim is true.
\[
 p_{\_}:T\to\operatorname{Generated}[A],\qquad
 e:A\to\operatorname{Generated}[S],\qquad s=e(p_{\_}(t).\mathrm{value}).
\]
Here $T$ denotes \texttt{Cohere\allowbreak Domain\allowbreak Target}, $A$ denotes \texttt{Cohere\allowbreak Scope\allowbreak Activation}, and $S$ denotes \texttt{Cohere\allowbreak Intake\allowbreak Disposition}. The displayed \texttt{.value} means a successfully obtained generated payload; it is neither a proof nor a rule for ignoring execution failure. No multi-input join occurs.

\paragraph{Source-prescribed conceptual work.}
The Yang--Mills bindings instantiate a domain-specific programme rather than a generic default. Its target-freezing procedure separates theory existence, Hamiltonian spectral exclusion, clustering, confinement, and area law; these are different conclusions with different burdens. The representation procedures then make explicit which observable, state, dynamics, and limit structures a proposed bridge preserves. The central question becomes which implication in the proposed programme remains unsupported, not how many steps have been completed. The Cohere workflow contract supplies a child-stage transition structure, while individual skill bindings provide methods within those stages; the binding list alone establishes no execution order. In this reading, intake identifies the precise programme whose claims will be investigated, and later stages can return a conditional result or an obstruction without pretending to settle its ultimate target.

As a constructed Yang--Mills example, a proposal starts from decay of a finite set of Wilson-loop observables and announces a full mass gap. Intake identifies the intended pure-gauge target and the distinct missing claims: observable completeness, physical-unit uniformity, continuum existence, and reconstruction. The useful outcome is a programme of separate obligations instead of accepting the announced implication. No decay or mass-gap result is asserted by this example.

The original policy sends ready intake to \texttt{cohere\_fast\_scope\_and\_ground} and stops blocked or stalled work. The native reasoning declaration only returns $S$; it does not invoke that next workflow or construct its grounding request. Its module-level input adapter remains separate from the policy arrow. Thus the workflow's conceptual role is delimiting the question on which later reasoning can operate, with an explicit native realization gap. This gap does not imply that the source procedures lack conceptual content.

% SPDX-FileCopyrightText: © 2026 Idiap Research Institute <contact@idiap.ch>
%
% SPDX-FileContributor: Andre Freitas <andre.freitas@idiap.ch>
% Neuro-symbolic AI Group
%
% SPDX-FileContributor: Danilo Gusicuma <danilo.gusicuma@idiap.ch>
% Neuro-symbolic AI Group
%
% SPDX-License-Identifier: MIT
% Generated from step-explanations.json
\paragraph{Step-by-step reconstruction.}
A proposed Yang--Mills programme claims that decay of finitely many Wilson-loop observables establishes the full mass gap.
The following is a source-guided illustrative reconstruction. Internal procedures are prescriptions; the native declarations do not establish their execution.
\begin{enumerate}[leftmargin=*,label=\arabic*.]
\item \textbf{Identify the programme} (\texttt{scope\_\allowbreak{}activation}).
\begin{description}[style=standard,labelindent=0pt,leftmargin=5.5em,labelwidth=5em,labelsep=0.5em,align=left,font=\normalfont\bfseries,itemsep=0.15em,parsep=0pt,topsep=0.3em]
\item[Purpose] Separate the scientific target from its support.
\item[Input] Activated domain target; intended claim
\item[Method] The native scope producer is unresolved. Target-freezing instructions distinguish existence, spectral exclusion, clustering, confinement and area law before selecting a programme.
\item[Check] Different endpoints require different premises and cannot be merged by a shared label.
\item[Output] Intended activation: pure-gauge programme; separate target burdens
\item[Next] Intake calibration consumes the activation; its existence does not establish any target theorem.
\item[Example] Split the announcement into finite-observable decay, observable completeness, continuum existence, reconstruction and spectral exclusion instead of accepting one as all the others.
\end{description}
\item \textbf{Judge whether inquiry can begin} (\texttt{intake\_\allowbreak{}disposition}).
\begin{description}[style=standard,labelindent=0pt,leftmargin=5.5em,labelwidth=5em,labelsep=0.5em,align=left,font=\normalfont\bfseries,itemsep=0.15em,parsep=0pt,topsep=0.3em]
\item[Purpose] Assess whether the scoped question is usable.
\item[Input] Scope activation; domain; target
\item[Method] Assess whether the target and its boundaries are definite enough to investigate, retaining missing claims as obligations.
\item[Check] Readiness concerns a usable question, not support for its ultimate conclusion.
\item[Output] Intake disposition: ready, blocked or stalled
\item[Next] Original policy routes ready intake to grounding; the native return constructs no grounding request.
\item[Example] Proceed only with the separate obligations visible; an unspecified meaning of mass gap requires target clarification before scientific grounding.
\end{description}
\end{enumerate}
\paragraph{Changed-input contrast.}
If intake is instead asked only for a finite-regulator observable estimate, continuum reconstruction is not part of that narrower result. The activation must preserve the selected target rather than silently exchanging the two questions.


\subsection{Scope and grounding: restricting the admissible claim}
\label{app:granular-cohere_fast_scope_and_ground}
\texttt{cohere\_fast\_scope\_and\_ground} makes domain, target, sources, and maximum claim scope explicit before discovery proceeds. Its unresolved projection transforms a request containing domain, target, and source references into a grounding record that also names a claim ceiling. The subsequent evidential assessment asks whether that bounded grounding is usable. These are distinct responsibilities: constructing a scope description and judging its adequacy. The original ERFL description gives no reasoning operator for the former, retaining \texttt{pre:\_}; naming it a projection does not implement the source-interpretation procedure prescribed by its linked skills.
\[
 p_{\_}:Q\to\operatorname{Generated}[G],\qquad
 e:G\to\operatorname{Generated}[S],\qquad s=e(p_{\_}(q).\mathrm{value}).
\]
Here $Q$ denotes \texttt{Domain\allowbreak Scope\allowbreak Grounding\allowbreak Request}, $G$ denotes \texttt{Domain\allowbreak Scope\allowbreak Grounding}, and $S$ denotes \texttt{Scope\allowbreak Grounding\allowbreak Disposition}. The generated payload projection is conditional on successful acquisition. The single relation, \texttt{Domain\allowbreak Scope\allowbreak Grounding\allowbreak Enables\allowbreak Disposition}, licenses assessment of that grounding as a dependency; it is not an inference from source presence to claim truth. There is no join.

\paragraph{Source-prescribed conceptual work.}
In the Yang--Mills instantiation, target freezing first specifies the pure gauge theory, compact simple gauge group, dimension, axiomatic endpoint, vacuum sector, and regulator, volume, cutoff, and coupling quantifiers. It partitions premises from targets: assuming a positive physical correlation scale would conceal the very estimate being sought. Literature grounding binds each imported claim to a locatable primary source and retains its conventions and dependencies. The conceptual product is an admissible implication problem with a precise scope, not merely a source list. A strong-coupling regulated statement, an abstract transfer theorem, and a continuum Hamiltonian claim cannot be interchanged. These procedures prescribe what the native grounding record is meant to summarize; its string fields do not establish that the interpretation or preserved scope is correct.

In a constructed Yang--Mills case, a cited estimate holds at fixed lattice spacing in a strong-coupling regime. Grounding retains that regime and asks which constants depend on volume and cutoff. It does not relabel the estimate as a continuum mass gap. The exact open question is whether a suitable physical decay rate survives the required limits under a justified comparison, rather than whether the source contains the words ``mass gap''.

Original accepted status routes to adaptive frontier formation; blocked status stops. Native code returns only the disposition. The next workflow's identically named \texttt{Domain\allowbreak Scope\allowbreak Grounding} is module-local and has a different field set; neither matching names nor the policy transition supplies an adapter. A handoff must preserve the intended target and claim ceiling explicitly. The unresolved grounding operation and this cross-module boundary remain visible rather than being replaced by an invented executable chain.

% SPDX-FileCopyrightText: © 2026 Idiap Research Institute <contact@idiap.ch>
%
% SPDX-FileContributor: Andre Freitas <andre.freitas@idiap.ch>
% Neuro-symbolic AI Group
%
% SPDX-FileContributor: Danilo Gusicuma <danilo.gusicuma@idiap.ch>
% Neuro-symbolic AI Group
%
% SPDX-License-Identifier: MIT
% Generated from step-explanations.json
\paragraph{Step-by-step reconstruction.}
The source packet contains a strong-coupling estimate at fixed lattice spacing. The intended endpoint is a physical continuum spectral statement.
The following is a source-guided illustrative reconstruction. Internal procedures are prescriptions; the native declarations do not establish their execution.
\begin{enumerate}[leftmargin=*,label=\arabic*.]
\item \textbf{Construct admissible grounding} (\texttt{scope\_\allowbreak{}grounding}).
\begin{description}[style=standard,labelindent=0pt,leftmargin=5.5em,labelwidth=5em,labelsep=0.5em,align=left,font=\normalfont\bfseries,itemsep=0.15em,parsep=0pt,topsep=0.3em]
\item[Purpose] Ground a bounded implication problem in sources.
\item[Input] Grounding request: domain; target; source references
\item[Method] The native producer is unresolved. Linked instructions freeze theory, gauge group, quantifiers and endpoint, then bind imported claims to locatable sources with their conventions.
\item[Check] Do not assume the desired physical scale or erase a source theorem's regulator regime.
\item[Output] Intended grounding: fixed-spacing support; explicit claim ceiling
\item[Next] Disposition assessment receives the grounding; nominal fields do not ensure faithful source interpretation.
\item[Example] Retain the estimate's strong-coupling and fixed-spacing assumptions and mark dependence on volume and cutoff as questions to inspect.
\end{description}
\item \textbf{Assess grounding adequacy} (\texttt{grounding\_\allowbreak{}disposition}).
\begin{description}[style=standard,labelindent=0pt,leftmargin=5.5em,labelwidth=5em,labelsep=0.5em,align=left,font=\normalfont\bfseries,itemsep=0.15em,parsep=0pt,topsep=0.3em]
\item[Purpose] Assess source support for further inquiry.
\item[Input] Grounding record; target; claim ceiling
\item[Method] Compare imported assumptions with the requested endpoint and identify the missing implication between their regimes.
\item[Check] Fixed-regulator support must not be relabeled as a continuum Hamiltonian conclusion.
\item[Output] Grounding disposition: usable bounded support or specified blocker
\item[Next] Original policy routes accepted grounding onward; module-local type names still require a preserving adapter.
\item[Example] Accept the source as a restricted starting point while keeping physical-rate uniformity and continuum passage open; block a packet that silently treats them as proved.
\end{description}
\end{enumerate}
\paragraph{Changed-input contrast.}
A source theorem uniform in the required limits would change the grounding burden, but only after its assumptions and theory identity were checked. Matching terminology is never sufficient transfer evidence.


\subsection{Adaptive frontier: identifying the next bounded question}
\label{app:granular-cohere_fast_adaptive_frontier}
\texttt{cohere\_fast\_adaptive\_frontier} organizes what to investigate next within an existing claim boundary. The first occurrence projects a grounding into a frontier reference with the intended preservation of domain and claim ceiling. Its string fields do not enforce equality with the input across this projection. Its native name, \texttt{Unresolved\allowbreak Domain\allowbreak Frontier\allowbreak Projection}, reflects the original \texttt{pre:\_}: the declaration specifies the frontier interface, while the linked procedures below prescribe how alternatives and representation burdens are investigated. Evidential readiness then judges whether the proposed frontier is suitable for the derivation lane.
\[
 p_{\_}:G\to\operatorname{Generated}[F],\qquad
 e:F\to\operatorname{Generated}[S],\qquad s=e(p_{\_}(g).\mathrm{value}).
\]
Here $G$ denotes \texttt{Domain\allowbreak Scope\allowbreak Grounding}, $F$ denotes \texttt{Bounded\allowbreak Domain\allowbreak Frontier}, and $S$ denotes \texttt{Frontier\allowbreak Readiness\allowbreak Disposition}, all local nominal types. Every output remains generated; \texttt{.value} denotes successful payload access. \texttt{Domain\allowbreak Frontier\allowbreak Enables\allowbreak Readiness\allowbreak Assessment} records the prerequisite for the assessment, with no additional scientific warrant and no multi-input join.

\paragraph{Source-prescribed conceptual work.}
The linked Yang--Mills procedures construct alternatives along three connected axes. A regulator atlas compares carriers, measures, boundary conditions, gauge invariance, locality, reflection positivity, and closure under blocking; it chooses a primary representation and an adversarial alternative by their missing theorem obligations. Observable-core construction asks whether the selected observables separate states and generate the physical sectors needed by a spectral argument, then seeks a hidden-sector countermodel. The representation calculus specifies observable pullbacks, state pushforwards, their compatibility, and observable-indexed approximation errors. It distinguishes equivalence from coarse-graining or a conjectural correspondence. Restricted-model calibration tests these proposals on massive, massless, finite-volume, and hidden-sector cases and revises definitions that produce false positives. This is the substantive frontier-building work compressed by the native projection: discovering which representation and observable choices expose a tractable missing implication, while preserving what could invalidate transfer.

A constructed diagnostic example uses a direct sum of a gapped sector and an unobserved gapless sector. An observable family restricted to the first sector can show decay while failing to detect the second. The frontier therefore changes from ``fit a faster decay curve'' to ``establish density of the observable-generated physical subspace, or report a sectorial claim.'' This is a calibration countermodel prescribed by the source procedures, not a claim that the Yang--Mills theory has such a decomposition.

The original policy sends accepted frontiers to \texttt{cohere\_fast\_derivation\_and\_proof} and stops quarantined ones. Native reasoning returns the readiness disposition, not the frontier needed by that next input. Exposing or reconstructing the frontier for a subsequent call requires a separate handoff. The native declaration itself does not implement an adaptive search loop: ``adaptive'' names the regime's purpose, while realization of its prescribed conceptual work as this native occurrence remains unresolved.

% SPDX-FileCopyrightText: © 2026 Idiap Research Institute <contact@idiap.ch>
%
% SPDX-FileContributor: Andre Freitas <andre.freitas@idiap.ch>
% Neuro-symbolic AI Group
%
% SPDX-FileContributor: Danilo Gusicuma <danilo.gusicuma@idiap.ch>
% Neuro-symbolic AI Group
%
% SPDX-License-Identifier: MIT
% Generated from step-explanations.json
\paragraph{Step-by-step reconstruction.}
In an abstract spectral calibration model, a visible gapped sector yields decaying correlations while an unobserved gapless sector coexists with it. This tests an inference, not existence of a gapless Yang--Mills theory. The frontier must expose the missing completeness burden.
The following is a source-guided illustrative reconstruction. Internal procedures are prescriptions; the native declarations do not establish their execution.
\begin{enumerate}[leftmargin=*,label=\arabic*.]
\item \textbf{Build a discriminating frontier} (\texttt{bounded\_\allowbreak{}frontier}).
\begin{description}[style=standard,labelindent=0pt,leftmargin=5.5em,labelwidth=5em,labelsep=0.5em,align=left,font=\normalfont\bfseries,itemsep=0.15em,parsep=0pt,topsep=0.3em]
\item[Purpose] Expose the missing implication through representation questions.
\item[Input] Grounded domain; fixed claim ceiling
\item[Method] The native producer is unresolved. Linked methods compare regulators, observable coverage and representation maps, then use restricted massive, massless and hidden-sector cases to test false positives.
\item[Check] A representation earns a role by preserving required structure and exposing a falsifiable burden.
\item[Output] Intended frontier: observable-density obligation or sectorial conclusion
\item[Next] Readiness assesses this frontier; the signature does not implement adaptive search.
\item[Example] The hidden-sector construction changes the next question from fitting decay to proving that the chosen observable class sees every relevant physical sector.
\end{description}
\item \textbf{Judge frontier readiness} (\texttt{frontier\_\allowbreak{}readiness}).
\begin{description}[style=standard,labelindent=0pt,leftmargin=5.5em,labelwidth=5em,labelsep=0.5em,align=left,font=\normalfont\bfseries,itemsep=0.15em,parsep=0pt,topsep=0.3em]
\item[Purpose] Assess whether the frontier supports bounded derivation.
\item[Input] Frontier; preserved domain and claim ceiling
\item[Method] Check whether the chosen representation names the objects, required map and missing theorem that derivation must address.
\item[Check] A promising representation without its transfer obligation is not an unconditional route.
\item[Output] Frontier disposition: ready or quarantined
\item[Next] The native return is a disposition, not the frontier needed by the next workflow.
\item[Example] A precise density obligation can enter derivation as an open task; a claim that visible decay already proves global exclusion must be quarantined.
\end{description}
\end{enumerate}
\paragraph{Changed-input contrast.}
Once the observable-generated physical subspace is known to be dense and the extension estimates are justified, the next frontier can move to remaining transfer and limit obligations. Repeating the same visible-sector fit does not supply density.


\subsection{Derivation and proof: exposing a conditional result}
\label{app:granular-cohere_fast_derivation_and_proof}
\texttt{cohere\_fast\_derivation\_and\_proof} demands a derivation packet whose unresolved obligations remain attached. Despite the workflow name, its original first occurrence is \texttt{pre:\_}, not a declared deduction operator. The native \texttt{Unresolved\allowbreak Conditional\allowbreak Derivation\allowbreak Projection} specifies a transformation from a bounded frontier into a packet containing a derivation reference, open-obligations reference, domain, and claim ceiling. Evidential assessment then asks whether that packet is ready for audit; it returns no proof verdict.
\[
 p_{\_}:F\to\operatorname{Generated}[D],\qquad
 e:D\to\operatorname{Generated}[S],\qquad s=e(p_{\_}(f).\mathrm{value}).
\]
Here $F$ denotes \texttt{Bounded\allowbreak Domain\allowbreak Frontier}, $D$ denotes \texttt{Conditional\allowbreak Derivation\allowbreak Packet}, and $S$ denotes \texttt{Derivation\allowbreak Readiness\allowbreak Disposition}. The successful \texttt{.value} adapter does not certify the derivation. \texttt{Derivation\allowbreak Packet\allowbreak Enables\allowbreak Readiness\allowbreak Assessment} marks a dependency between packet and assessment, without converting readiness into entailment. There is no join of independent evidence here.

\paragraph{Source-prescribed conceptual work.}
The Yang--Mills proof bindings go substantially beyond packaging a conditional formula. Adaptive block domination fixes the exposure filtration and retained boundary information, bounds normalized probability on a good chart, and charges the complementary configurations through a multi-link estimate. Heat-kernel prefactors, chart mass, Haar Jacobians, overlap, and remainder oscillation enter the normalized comparison; a pointwise action estimate alone does not supply it. A state-weighted alternative proves a positive-kernel drift inequality $QV\leq q_VV$, controls the initial weight moment, and derives iteration before checking a strict cluster threshold. Retained-state sufficiency, positive-mass histories, and uniformity in volume and boundary conditions are independent obligations. Finite-block and enclosure candidates make parts of these burdens checkable; exact finite checks do not discharge full-measure, repeated-RG, continuum, or reconstruction hypotheses. The prescribed broader programme separates basin entry, controlled iteration, physical-unit decay, positivity, observable density, and continuum passage. These are implication dependencies to prove, not guaranteed successes or an execution order inferred from the skill list.

For a constructed finite diagnostic, the law $P(10)=P(01)=1/2$ has vanishing two-site intersection probability, yet $P(X_2=1\mid X_1=0)=1$. A small intersection moment therefore cannot replace the history-uniform conditional ceiling required by the domination argument. The next proof task is to identify an exposure state and prove its normalized conditional bound, not to declare domination from that moment. This source-prescribed counterexample tests a proposed proof mechanism; it proves no Yang--Mills, continuum, or Hamiltonian-gap conclusion.

Original accepted status routes to independent audit; quarantine stops. Native reasoning returns only readiness, whereas audit requires the actual conditional packet. That handoff is not executable merely because the workflow names are adjacent. Both the missing derivation producer and the preservation of open obligations must remain explicit when describing the larger chain.

% SPDX-FileCopyrightText: © 2026 Idiap Research Institute <contact@idiap.ch>
%
% SPDX-FileContributor: Andre Freitas <andre.freitas@idiap.ch>
% Neuro-symbolic AI Group
%
% SPDX-FileContributor: Danilo Gusicuma <danilo.gusicuma@idiap.ch>
% Neuro-symbolic AI Group
%
% SPDX-License-Identifier: MIT
% Generated from step-explanations.json
\paragraph{Step-by-step reconstruction.}
A proposed domination argument treats a small joint event probability as a history-uniform conditional ceiling. Test the finite law assigning probability one half each to outcomes 10 and 01.
The following is a source-guided illustrative reconstruction. Internal procedures are prescriptions; the native declarations do not establish their execution.
\begin{enumerate}[leftmargin=*,label=\arabic*.]
\item \textbf{Construct the conditional derivation} (\texttt{conditional\_\allowbreak{}derivation}).
\begin{description}[style=standard,labelindent=0pt,leftmargin=5.5em,labelwidth=5em,labelsep=0.5em,align=left,font=\normalfont\bfseries,itemsep=0.15em,parsep=0pt,topsep=0.3em]
\item[Purpose] Specify the mechanism’s required inequality and assumptions.
\item[Input] Bounded frontier; target; claim ceiling
\item[Method] The native producer is unresolved. Linked domination methods specify exposure history, retained boundary information and normalized conditional bounds; finite diagnostics attack the required implication.
\item[Check] Joint moments cannot replace history-uniform bounds; normalization and positive-probability conditioning must be justified.
\item[Output] Intended derivation packet: failed inference; open replacement obligation
\item[Next] Readiness assesses the packet; the declared projection is not a proof-producing algorithm.
\item[Example] For P(10)=P(01)=1/2, the two-site intersection probability is zero, but P(X2=1 given X1=0)=1. Thus the suggested conditional ceiling does not follow.
\end{description}
\item \textbf{Assess readiness for audit} (\texttt{derivation\_\allowbreak{}readiness}).
\begin{description}[style=standard,labelindent=0pt,leftmargin=5.5em,labelwidth=5em,labelsep=0.5em,align=left,font=\normalfont\bfseries,itemsep=0.15em,parsep=0pt,topsep=0.3em]
\item[Purpose] Assess whether the derivation is auditable.
\item[Input] Conditional packet; open-obligation reference
\item[Method] Check that the proposed inequality, retained state, assumptions and failure cases remain explicit rather than hidden by a certificate label.
\item[Check] A finite check cannot discharge full-measure, repeated-iteration or continuum hypotheses.
\item[Output] Derivation disposition: ready for scoped audit or quarantined
\item[Next] Original policy sends accepted packets to audit; native return omits the packet adapter.
\item[Example] Send the counterexample and the requirement for a normalized history-uniform bound to audit; withhold any domination conclusion inferred only from the zero intersection moment.
\end{description}
\end{enumerate}
\paragraph{Changed-input contrast.}
A proved normalized conditional bound for every admissible positive-probability exposure history would answer this particular objection. Small unconditional intersection moments alone do not, and even that bound would not settle continuum reconstruction.


\subsection{Independent audit: challenging a claim without losing it}
\label{app:granular-cohere_fast_independent_audit}
\texttt{cohere\_fast\_independent\_audit} has two declared reasoning mechanisms. Refutation examines a conditional derivation for a defeater and returns \texttt{Refuted}, \texttt{Not\allowbreak Refuted\allowbreak Within\allowbreak Scope}, or \texttt{Inconclusive}. Calibration then receives both the original derivation and that assessment. This join matters because the same adverse finding can refute an overstrong conclusion while leaving a narrower conditional implication intact. The relation \texttt{Calibrates\allowbreak Independent\allowbreak Audit\allowbreak Disposition} licenses use of the audit finding to assess the retained packet; it supplies no theorem independently of the finding.
\[
 r:D\to\operatorname{Generated}[V],\quad
 e:J\to\operatorname{Generated}[S].
\]
\[
 s=e\bigl(\{\mathrm{derivation}=d,\mathrm{refutation\_assessment}=r(d).\mathrm{value}\}\bigr).
\]
Here $D$ denotes \texttt{Conditional\allowbreak Derivation\allowbreak Packet}, $V$ denotes \texttt{Independent\allowbreak Refutation\allowbreak Assessment}, $J$ denotes \texttt{Independent\allowbreak Audit\allowbreak Disposition\allowbreak Input}, and $S$ denotes \texttt{Independent\allowbreak Audit\allowbreak Disposition}. Generated assessment access is a successful-result view. $J$ is a nominal record, not a native dependent pair enforcing the assessment's exact claim index.

\paragraph{Source-prescribed conceptual work.}
The linked Yang--Mills audit expands every cutoff, volume, distance, observable, and coupling quantifier, then seeks the first unproved implication in the complete proposed argument. It checks whether a bridge encodes its desired conclusion in its definition, whether constants remain uniform in physical units, and whether a representation change loses positivity, locality, or observable completeness. Limit interchanges require their own justification: a favorable order of limits is not automatically licensed. The audit distinguishes stochastic mixing, finite transfer spectra, sectorial decay, and a continuum Hamiltonian gap. Countermodels and the finite checkable core are used to attack particular steps rather than to cast a generic confidence vote. Its output is the strongest defensible statement together with the exact missing bridge and required downgrade language. This source-prescribed adversarial programme explains what the native refutation assessment summarizes; its name alone supplies no empirical independence guarantee.

In a constructed audit, a proposed spectral argument supplies temporal decay only for a finite observable family. The hidden-sector countermodel challenges the step from those correlations to the full spectrum. Refutation targets the missing density implication; calibration can retain a conditional or sectorial result while withholding the full claim. This differs from rejecting the decay calculation itself. Bounded non-refutation of the calculation cannot fill the missing density proof.

Original accepted status routes to manuscript release and quarantined status stops. Native code returns $S$ only. Release instead expects an \texttt{Independent\allowbreak Audit\allowbreak Packet} containing an audit reference and claim ceiling, so a separate construction must retain the reviewed substance. No unresolved \texttt{\_} occurrence appears inside this audit, but the cross-workflow packet adapter remains outside its signature.

% SPDX-FileCopyrightText: © 2026 Idiap Research Institute <contact@idiap.ch>
%
% SPDX-FileContributor: Andre Freitas <andre.freitas@idiap.ch>
% Neuro-symbolic AI Group
%
% SPDX-FileContributor: Danilo Gusicuma <danilo.gusicuma@idiap.ch>
% Neuro-symbolic AI Group
%
% SPDX-License-Identifier: MIT
% Generated from step-explanations.json
\paragraph{Step-by-step reconstruction.}
A conditional spectral argument relies on decay for finitely many observables and leaves density of their generated physical sector unproved.
The following is a source-guided illustrative reconstruction. Internal procedures are prescriptions; the native declarations do not establish their execution.
\begin{enumerate}[leftmargin=*,label=\arabic*.]
\item \textbf{Attack the complete implication} (\texttt{refutation\_\allowbreak{}assessment}).
\begin{description}[style=standard,labelindent=0pt,leftmargin=5.5em,labelwidth=5em,labelsep=0.5em,align=left,font=\normalfont\bfseries,itemsep=0.15em,parsep=0pt,topsep=0.3em]
\item[Purpose] Locate the calculation-to-endpoint inference gap.
\item[Input] Conditional derivation; quantifiers; open obligations
\item[Method] Expand volume, cutoff, distance and observable quantifiers; test hidden sectors, limit interchanges and definitions that encode the desired conclusion.
\item[Check] A defect must target a particular inference; bounded non-refutation is not full proof.
\item[Output] Refutation assessment: missing density implication or another scoped defeater
\item[Next] Calibration receives this assessment with the original derivation; the name independent proves no empirical independence.
\item[Example] As an abstract spectral countermodel, adjoin an orthogonal sector with positive energies approaching zero, invisible to all tested observables. Their decay survives, so the general global-gap inference fails; this does not construct a Yang--Mills counterexample.
\end{description}
\item \textbf{Retain the defensible claim} (\texttt{audit\_\allowbreak{}disposition}).
\begin{description}[style=standard,labelindent=0pt,leftmargin=5.5em,labelwidth=5em,labelsep=0.5em,align=left,font=\normalfont\bfseries,itemsep=0.15em,parsep=0pt,topsep=0.3em]
\item[Purpose] Calibrate the claim while preserving its valid part.
\item[Input] Original derivation + its refutation assessment
\item[Method] Compare the audited scope with the failure, identify the strongest retained statement and state the required downgrade.
\item[Check] Do not confuse refuting an overstrong extension with refuting its narrower conditional calculation.
\item[Output] Audit disposition: retained sectorial support; quarantined full-spectrum inference
\item[Next] Release requires a separate audit packet preserving this content; native disposition alone is insufficient.
\item[Example] Keep the visible-sector decay statement and make density an explicit missing bridge, preventing a favorable local audit from certifying a global mass gap.
\end{description}
\end{enumerate}
\paragraph{Changed-input contrast.}
If the packet states only visible-sector decay, the hidden-sector example does not refute that restricted calculation. It defeats the stronger global inference, so calibration must preserve the distinction.


\subsection{Manuscript release: preserving authority while presenting a result}
\label{app:granular-cohere_fast_manuscript_release}
\texttt{cohere\_fast\_manuscript\_release} asks whether a presentation stays within the authority of its source audit. Its first occurrence is intended to construct a presentation reference tied to the audit and its claim ceiling. String references alone do not enforce that cross-projection preservation. The original \texttt{pre:\_} and native \texttt{Unresolved\allowbreak Claim\allowbreak Bounded\allowbreak Presentation\allowbreak Projection} leave the native realization unresolved; the linked procedures prescribe the paper-construction work described below. The second occurrence calibrates release disposition from that presentation. It assesses an exposition, not a new scientific argument that strengthens its source claims.
\[
 p_{\_}:A\to\operatorname{Generated}[P],\qquad
 e:P\to\operatorname{Generated}[S],\qquad s=e(p_{\_}(a).\mathrm{value}).
\]
Here $A$ denotes \texttt{Independent\allowbreak Audit\allowbreak Packet}, $P$ denotes \texttt{Claim\allowbreak Bounded\allowbreak Presentation}, and $S$ denotes \texttt{Manuscript\allowbreak Release\allowbreak Disposition}. The payload adapter remains conditional on successful generation. The single relation \texttt{Claim\allowbreak Bounded\allowbreak Presentation\allowbreak Enables\allowbreak Release} expresses the prerequisite for the release assessment; it adds no evidence for the presented scientific proposition. There is no multi-input join: the source audit reference and ceiling must already be retained in $P$.

\paragraph{Source-prescribed conceptual work.}
The domain-private writing bindings divide presentation into organizing accepted results and explicit open obligations, assembling proof details without upgrading their status, and rendering the manuscript from the accepted evidence packets. Each procedure retains conventions, dependencies, attempted falsifiers, failed routes, and conditional branches. The resulting argument must show which theorem uses which hypothesis and where a finite certificate stops short of its intended application. This is more than checking that a file compiles: a reader must be able to distinguish established mathematical content from an unproved bridge. The procedures prescribe these construction responsibilities without specifying a complete automatic prose-generation algorithm. They apply to the explicitly activated Yang--Mills instance; they are not a default interpretation of every scientific paper. Their ordering is not inferred merely from appearing together in a binding list.

As a constructed paper example, a finite-block certificate checks a normalized bound for declared boundary data, while a volume-uniform supremum remains open. The narrative states the finite result; the proof appendix exposes the exact parameters and checked inequality; the open supremum stays a named obligation. An abstract claiming uniform control of all blocks would change the result even if the finite computation were reproduced perfectly. Faithful construction makes the useful certificate intelligible without pretending that the broader implication has been proved.

Original accepted status accepts the presentation; stopped status stops. Native reasoning only returns that disposition and does not publish a paper or execute an acceptance branch. The manifest leaves materialization blocked on the unresolved presentation projection. Thus the workflow describes a scientific communication boundary with an explicit native realization gap, rather than an automatic publication mechanism or a further proof step.

% SPDX-FileCopyrightText: © 2026 Idiap Research Institute <contact@idiap.ch>
%
% SPDX-FileContributor: Andre Freitas <andre.freitas@idiap.ch>
% Neuro-symbolic AI Group
%
% SPDX-FileContributor: Danilo Gusicuma <danilo.gusicuma@idiap.ch>
% Neuro-symbolic AI Group
%
% SPDX-License-Identifier: MIT
% Generated from step-explanations.json
\paragraph{Step-by-step reconstruction.}
An audit supports a finite-block certificate for declared boundary data but leaves a volume-uniform supremum estimate open.
The following is a source-guided illustrative reconstruction. Internal procedures are prescriptions; the native declarations do not establish their execution.
\begin{enumerate}[leftmargin=*,label=\arabic*.]
\item \textbf{Construct a faithful presentation} (\texttt{claim\_\allowbreak{}bounded\_\allowbreak{}presentation}).
\begin{description}[style=standard,labelindent=0pt,leftmargin=5.5em,labelwidth=5em,labelsep=0.5em,align=left,font=\normalfont\bfseries,itemsep=0.15em,parsep=0pt,topsep=0.3em]
\item[Purpose] Present the finite result within its supported scope.
\item[Input] Independent audit reference; claim ceiling
\item[Method] The native presentation producer is unresolved. Linked writers organize accepted results, expose proof dependencies and retain failed routes and open obligations while rendering the argument.
\item[Check] Each theorem must state the hypotheses it uses and where the finite certificate stops.
\item[Output] Intended presentation: finite certificate; unresolved uniform supremum
\item[Next] Release assessment consumes the presentation; string references do not enforce preservation automatically.
\item[Example] State exactly which boundary data the certificate checks, show the finite inequality, and list the supremum estimate as a separate open application burden.
\end{description}
\item \textbf{Assess release fidelity} (\texttt{release\_\allowbreak{}disposition}).
\begin{description}[style=standard,labelindent=0pt,leftmargin=5.5em,labelwidth=5em,labelsep=0.5em,align=left,font=\normalfont\bfseries,itemsep=0.15em,parsep=0pt,topsep=0.3em]
\item[Purpose] Assess manuscript fidelity to the audit.
\item[Input] Presentation; retained audit reference and claim ceiling
\item[Method] Compare headline, theorem and application claims against the accepted finite statement and remaining obligations.
\item[Check] Readability and compilation cannot strengthen the source claim.
\item[Output] Release disposition: supportable presentation or stopped overstatement
\item[Next] Original acceptance is a policy; native return neither publishes nor proves the result.
\item[Example] Release only the finite-result narrative with the uniformity gap visible; stop an abstract announcing all-block control from that same certificate.
\end{description}
\end{enumerate}
\paragraph{Changed-input contrast.}
If the abstract changes the finite claim into uniform control over all blocks, the presentation must fail its claim boundary even when every finite calculation and LaTeX build succeeds.


\subsection{Result elaboration: extending a manuscript within its evidence}
\label{app:granular-manuscript_result_elaboration}
\texttt{manuscript\_result\_elaboration} combines four occurrences. Initial calibration extracts the scope of a draft; abduction proposes an extension within it; refutation tests that extension; final calibration joins the same extension with its validity assessment. The three relations distinguish their licenses: \texttt{Constrains\allowbreak Manuscript\allowbreak Extension} carries the scope restriction, \texttt{Tests\allowbreak Manuscript\allowbreak Extension} exposes the proposal to criticism, and \texttt{Calibrates\allowbreak Manuscript\allowbreak Readiness} supports a readiness judgment from proposal plus verdict.
\[
 e_0:P\to\operatorname{Generated}[Q],\quad a:Q\to\operatorname{Generated}[H],
\]
\[
 r:H\to\operatorname{Generated}[V],\quad e_1:J\to\operatorname{Generated}[S].
\]
\begingroup\raggedright
Here $P$ denotes \texttt{Paper\allowbreak Draft}, $Q$ denotes \texttt{Manuscript\allowbreak Result\allowbreak Scope}, $H$ denotes \texttt{Manuscript\allowbreak Extension\allowbreak Hypothesis}, $V$ denotes \texttt{Extension\allowbreak Validity\allowbreak Assessment}, $J$ denotes \texttt{Manuscript\allowbreak Readiness\allowbreak Input}, and $S$ denotes \texttt{Manuscript\allowbreak Readiness\allowbreak Assessment}. For an input draft $p$, write $q=e_0(p)$, $z=a(q.\mathrm{value})$, $h=z.\mathrm{value}$, and $v=r(h)$ after successful generation. The final call to $e_1$ consumes $J$ with \texttt{extension\_hypothesis}$=h$ and \texttt{validity\_assessment}$=v.\mathrm{value}$. Final $S$ contains disposition, integrated draft, expansion validation, and readiness decision. These nominal records do not automatically enforce scientific scope or dependent identity.\par\endgroup

\paragraph{Source-prescribed conceptual work.}
The result-elaboration skills first recover the draft's smallest honest contribution and classify each expansion marker by its mathematical role. A missing coefficient, finite-object definition, transfer theorem, uniform estimate, and explanatory paragraph require different work. Open-module decomposition separates exact algebra or finite matrix checking from generative construction of estimates and proof strategies. The planner orders needs by genuine dependence: coefficients precede sign tests, and transfer certificates precede stronger application claims. Finite-object formalization specifies the carrier, indices, equivalence or support rule, quantitative measure, worked non-example, and the lemma that may consume the object. Proof-packet derivation then states each inference from its definitions and hypotheses, localizing a missing lemma or countermodel when it fails. Prior-art and disqualifier review test novelty and scope; narrative integration explains the result in the argument. Thus the four native occurrences compress an evidence-sensitive programme of mathematical elaboration, not a single generic request to expand prose.

In a constructed finite-result manuscript, an open module requests a sign certificate for a parameter-dependent matrix. Elaboration first defines the matrix, parameter range, and error terms; coefficient extraction then supplies the quantities needed for an exact sign test. A proof packet explains how that test enters the stated lemma. If an application additionally requires a geometric transfer estimate, finite sign success leaves that estimate open. The manuscript receives a checkable finite contribution and a precise boundary, rather than a paragraph asserting that the application now follows.

Original \texttt{proof\_required} routes to theorem delivery, \texttt{paper\_ready} accepts, and stalled reasoning stops. Native code returns the readiness record; comments describe a non-executable adapter to \texttt{Proof\allowbreak Queue\allowbreak Item} and a separate acceptance gate. There is no unresolved \texttt{\_} occurrence, but generation, integration, and a ready label alone do not prove an extension or execute publication.

% SPDX-FileCopyrightText: © 2026 Idiap Research Institute <contact@idiap.ch>
%
% SPDX-FileContributor: Andre Freitas <andre.freitas@idiap.ch>
% Neuro-symbolic AI Group
%
% SPDX-FileContributor: Danilo Gusicuma <danilo.gusicuma@idiap.ch>
% Neuro-symbolic AI Group
%
% SPDX-License-Identifier: MIT
% Generated from step-explanations.json
\paragraph{Step-by-step reconstruction.}
A finite-result draft requests a sign certificate for the matrix with diagonal entries 1 and off-diagonal entries t, for absolute t at most one half. A geometric application still needs a separate transfer estimate.
The following is a source-guided illustrative reconstruction. Internal procedures are prescriptions; the native declarations do not establish their execution.
\begin{enumerate}[leftmargin=*,label=\arabic*.]
\item \textbf{Recover the actual expansion need} (\texttt{e0}).
\begin{description}[style=standard,labelindent=0pt,leftmargin=5.5em,labelwidth=5em,labelsep=0.5em,align=left,font=\normalfont\bfseries,itemsep=0.15em,parsep=0pt,topsep=0.3em]
\item[Purpose] Distinguish mathematical gaps from prose gaps.
\item[Input] Draft; stated contribution; expansion markers
\item[Method] Classify each request as a definition, coefficient, finite check, transfer theorem or explanatory paragraph and recover the smallest honest scope.
\item[Check] Do not treat all expansion markers as requests to assert stronger results.
\item[Output] Expansion scope: finite matrix certificate; separate application obligation
\item[Next] Extension formation consumes this scope; unresolved geometric work remains distinct.
\item[Example] Specify the two-by-two symmetric matrix, the range absolute t at most one half, and the still-unproved geometric transfer estimate.
\end{description}
\item \textbf{Propose a bounded extension} (\texttt{a0}).
\begin{description}[style=standard,labelindent=0pt,leftmargin=5.5em,labelwidth=5em,labelsep=0.5em,align=left,font=\normalfont\bfseries,itemsep=0.15em,parsep=0pt,topsep=0.3em]
\item[Purpose] Construct the missing finite certificate.
\item[Input] Matrix scope; parameter range; missing sign certificate
\item[Method] Define the carrier and entries, order coefficient extraction before the sign test, and state the lemma that will consume it.
\item[Check] A finite extension must expose its quantitative range and avoid assuming the application.
\item[Output] Extension hypothesis: exact eigenvalue-based positivity certificate
\item[Next] Refutation tests this proposed extension and its parameter boundary.
\item[Example] For the matrix with diagonal 1 and off-diagonal t, propose eigenvalues 1+t and 1-t and hence a lower bound of one half on the stated range.
\end{description}
\item \textbf{Test the extension's validity} (\texttt{r0}).
\begin{description}[style=standard,labelindent=0pt,leftmargin=5.5em,labelwidth=5em,labelsep=0.5em,align=left,font=\normalfont\bfseries,itemsep=0.15em,parsep=0pt,topsep=0.3em]
\item[Purpose] Test the finite result and its application boundary.
\item[Input] Proposed matrix certificate; range; intended use
\item[Method] Check the algebra and boundary values, seek a non-example, and localize any absent transfer lemma.
\item[Check] Finite sign success does not discharge a geometric comparison hypothesis.
\item[Output] Validity assessment: scoped matrix calculation; open transfer debt
\item[Next] Readiness joins this assessment with the same extension hypothesis.
\item[Example] At t=1/2 the eigenvalues are 3/2 and 1/2; at t=2 they are 3 and -1. The range matters, and neither calculation supplies geometric transfer.
\end{description}
\item \textbf{Integrate with its evidence boundary} (\texttt{e1}).
\begin{description}[style=standard,labelindent=0pt,leftmargin=5.5em,labelwidth=5em,labelsep=0.5em,align=left,font=\normalfont\bfseries,itemsep=0.15em,parsep=0pt,topsep=0.3em]
\item[Purpose] Assess inclusion while preserving proof obligations.
\item[Input] Extension hypothesis + its validity assessment
\item[Method] Explain how the finite lemma fits the argument, retain the counterexample and separate paper readiness from unproved application claims.
\item[Check] An integrated paragraph cannot turn a missing bridge into an accepted theorem.
\item[Output] Finite certificate; expansion validation; readiness decision
\item[Next] Original proof-required routing needs a ProofQueueItem adapter; native return does not execute it.
\item[Example] Add the exact finite positivity result and its range to the draft, while sending the geometric transfer estimate to further proof work instead of claiming the application.
\end{description}
\end{enumerate}
\paragraph{Changed-input contrast.}
If t is allowed to exceed 1, the smaller eigenvalue becomes negative and the proposed positivity extension fails. If the finite range is retained, positivity still cannot establish the missing geometric transfer.



